\chapter{自感、互感与磁路}
电感计算需要确定闭合电流路径及其磁链。各匝链接相同磁通时，磁链等于匝数乘单匝磁通；有限截面导体中的不同电流丝则可能链接不同磁通。计算机械力还需明确控制条件：电源保持电流不变，或无阻回路保持磁链不变。能量函数及求导方式由这些条件决定。

以下从场积分建立电感，再讨论耦合回路、铁芯磁路与多绕组变压器。除明确讨论导线内部或不同介质外，周围介质取真空；磁路采用线性、无磁滞模型。



\section{磁通、磁链与电感的定义}\label{sec:i2-linkage}
多匝线圈各匝的磁通未必相同，回路方程使用它们的有向总和，即磁链。电流正向和绕向一经规定，线性系统中的磁链便与各回路电流成线性关系。

\subsection{从一匝到多匝}
单匝闭路的外加电流为$I$，在线性静磁近似下，各点磁场均与$I$成正比。因此链接回路的磁通可写成$\Phi=LI$。对多匝线圈，应把各匝磁通相加，定义磁链
\[\lambda=\sum_{k=1}^N\Phi_k=LI.
\]
当每匝都链接同一磁通时，才可简化成$\lambda=N\Phi$。若匝数密度为$n$、线圈长度为$H$，总匝数是$N=nH$；将$n$误当成$N$会丢失一个长度因子。

线圈几何也可变化，此时感应电动势为
\[\boxed{\mathcal E_{\rm self}=-\frac{\dd(LI)}{\dd t}
=-L\dot I-I\dot L}.
\]
第一项来自电流变化，第二项来自几何或材料参数变化。只有$L$恒定时才可写成$-L\dot I$。多回路推广为$\lambda_i=\sum_jL_{ij}I_j$，每个$L_{ij}$都与所选回路正向相联系。

\subsection{有限截面导体的链接权重}
实心导体可分成许多平行电流丝。穿过导体内部的一小段磁通只链接外侧的一部分电流，而不是链接整个$I$。因此任意选一条穿过截面的切线，再把其磁通除以总电流，并不一定给出电感。

若用电流丝的链接描述，总磁链需按$\dd I/I$加权平均；若用磁能描述，则可直接由$W_B=LI^2/2$反求$L$。后一个办法自动包含各电流丝的自感和互感，更适合处理相邻的有限半径导线。对磁矢势，闭合稳恒电流还可使用$W_B=\frac12\int\vct j\cdot\vct A\,\dd V$；完整回流路径及相应边界项必须保留。
这一能量式可由场积分推出。在无额外永久磁化的线性静磁问题中，$\nabla\times\vct H=\vct j$、$\nabla\times\vct A=\vct B$，故
\[\nabla\cdot(\vct A\times\vct H)
=\vct H\cdot(\nabla\times\vct A)
-\vct A\cdot(\nabla\times\vct H)
=\vct B\cdot\vct H-\vct A\cdot\vct j.
\]
积分并乘$1/2$，得到
\[\frac12\int_V\vct B\cdot\vct H\dd V
=\frac12\int_V\vct j\cdot\vct A\dd V
+\frac12\oint_{\partial V}(\vct A\times\vct H)\cdot\vct n\dd S.
\]
对局域闭合电流，将积分面推到无穷远，表面项消失，才得到前述全空间等式。无限导线或有限截取区域须保留边界项；内部电感的直接场积分将在习题~\ref{ex:i2-h3}中给出。



\section{磁能、磁共能与外源功}\label{sec:i2-energy}
固定电流与固定磁链的形变具有不同的外源功。广义力由完整能量收支得到，采用磁能或磁共能时，偏导中保持的量必须符合题设。

\subsection{线性回路的建立过程}
固定几何，把电流从零准静态建立到$I$。电源增加的电磁功为$I\dd\lambda$，故
\[W_B=\int_0^\lambda I(\lambda')\dd\lambda'
=\frac1L\int_0^\lambda\lambda'\dd\lambda'
=\frac1L\left[\frac{\lambda'^2}{2}\right]_0^\lambda
=\frac12LI^2.
\]
线性介质中，它等于$\frac12\int\vct B\cdot\vct H\,\dd V$；真空中密度为$B^2/(2\mu_0)$。场随电流逐渐建立，积分因而给出系数$1/2$。用最终电流$I$乘全部磁链增长会把建立过程的功多算一倍。

若电感也随机械坐标$x$改变，
\[\dd\lambda=L\dd I+I L'(x)\dd x.
\]
电源功与磁能变化分别为
\[\dd W_{\rm src}=LI\dd I+I^2L'\dd x,\qquad
\dd W_B=LI\dd I+\frac12I^2L'\dd x.
\]
两者之差才是磁场对机械部分所作的功：
\[\boxed{F_x=\frac12I^2L'(x)}.
\]
在固定电流的虚位移中，磁能增加并不意味着磁力必然阻碍运动；维持电流的电源同时输入了两倍于这部分磁能增量的功。

\subsection{用磁链作变量}
无电阻、没有另接电源的闭路满足磁链守恒。用$\lambda$表示能量，
\[W_B(\lambda,x)=\frac{\lambda^2}{2L(x)}.
\]
此时机械力应按固定$\lambda$求负导数，
\[F_x=-\left.\partial_xW_B\right|_\lambda
=-\frac{\lambda^2}{2}\left[-L(x)^{-2}L'(x)\right]
=\frac{\lambda^2L'}{2L^2}=\frac12I^2L'.
\]
瞬时力与固定电流公式一致，但沿运动过程$I=\lambda/L(x)$会改变。两种控制条件不能共用同一条$I(x)$。

多回路的线性能量为$W_B=\frac12\vct I^{\mathsf T}\mathsf L\vct I$。若固定电流，使用磁共能$W'_B=\vct I^{\mathsf T}\vct\lambda-W_B$的正导数；在线性系统中$W'_B$与$W_B$数值相同，却具有不同的自然变量。固定磁链时使用$W_B=\frac12\vct\lambda^{\mathsf T}\mathsf L^{-1}\vct\lambda$的负导数。材料非线性时二者一般不相等；有磁滞时还不能只用一个单值能量函数描述全部过程。


\begin{example}[长螺线管与端部修正]\label{ex:i2-e1}
半径$a$、长度$H$的密绕空气螺线管，每单位长度$n$匝，通电流$I$。先用磁链与磁能两种方法求$H\gg a$时的自感，再求有限长度在$a/H$一阶上的修正，说明近似条件。 见图~\ref{fig:i2-e1-q}。
\end{example}

\begin{figure}[H]\centering
\begin{tikzpicture}[>=Stealth,line cap=round,line join=round]
\draw[curve](-2,-.75)--(2,-.75);\draw[curve](-2,.75)--(2,.75);
\draw[curve](-2,0)ellipse(.25 and .75);\draw[curve](2,0)ellipse(.25 and .75);
\foreach\x in{-1.7,-1.2,-.7,-.2,.3,.8,1.3,1.8}{\draw(\x,0)ellipse(.23 and .75);}
\draw[axis](-2.7,0)--(2.7,0)node[right]{$z$};\draw[->](-.8,.25)--(.8,.25)node[midway,above]{$\vct B$};
\draw[<->](-2,-1.15)--(2,-1.15)node[midway,below]{$H$};\node[above]at(0,.9){$n,\ I$};\draw[<->](2.5,0)--(2.5,.75)node[midway,right]{$a$};
\end{tikzpicture}
\caption{轴向均匀场只适用于长管中部，端部场需要另计。}\label{fig:i2-e1-q}
\end{figure}
\begin{solution}
图~\ref{fig:i2-e1-q}中，取远离两端的长螺线管。Ampere定律给$B\simeq\mu_0nI$，各匝链接磁通近似为$\pi a^2B$。总磁链
\[\lambda=(nH)\pi a^2(\mu_0nI),\qquad
\boxed{L_\infty=\mu_0n^2\pi a^2H}.
\]
磁能法给$W_B\simeq B^2(\pi a^2H)/(2\mu_0)=L_\infty I^2/2$，与链接法相同。这里$n$固定表示规定了绕组的匝数密度，不是已规定机械伸长时$n$不变。


\begin{figure}[H]\centering
\begin{tikzpicture}[>=Stealth,line cap=round,line join=round]
\begin{scope}[xshift=-2.5cm]\draw[curve](-1.6,-.65)--(1.6,-.65);\draw[curve](-1.6,.65)--(1.6,.65);
\draw[guide](-1.6,0)ellipse(.25 and .65);\draw[guide](1.6,0)ellipse(.25 and .65);\draw[->](-.8,0)--(.8,0)node[midway,above]{$B\simeq\mu_0nI$};
\draw[->](-1.3,.35)..controls(-2,1.2)and(-2,-1.2)..(-1.3,-.35);\draw[->](1.3,-.35)..controls(2,-1.2)and(2,1.2)..(1.3,.35);
\node[below]at(0,-1.15){端部漏磁降低链接};\end{scope}
\begin{scope}[xshift=2.7cm,x=.55cm]\draw[axis](3,0)--(9,0)node[right]{$H/a$};\draw[axis](3,0)--(3,1.8)node[above]{$L/L_\infty$};
\draw[curve,domain=3:8.7]plot(\x,{1-8/(3*pi*\x)});\draw[guide](3,1)--(8.7,1);\node[left]at(3,1){$1$};\end{scope}
\end{tikzpicture}
\caption{有限长度使自感低于无限长局部模型的结果。}\label{fig:i2-e1-s}
\end{figure}

图~\ref{fig:i2-e1-s}左图的端部漏磁给出首阶修正。将密绕线圈看作圆柱面电流片，对两匝的相对角$\phi$与轴向距离$u$积分。两圆周元的点积给$a^2\cos\phi$，横向距离为$d=2a\sin(\phi/2)$。Neumann积分整理为
\[L=\frac{\mu_0n^2a^2}{2}\int_0^{2\pi}\cos\phi\,K_H(d)\dd\phi,
\]
其中轴向积分为
\[\begin{aligned}
K_H(d)&=\int_{-H}^H\frac{H-|u|}{\sqrt{u^2+d^2}}\dd u\\
&=2H\int_0^H\frac{\dd u}{\sqrt{u^2+d^2}}
-2\int_0^H\frac{u\dd u}{\sqrt{u^2+d^2}}\\
&=2\left[H\arsinh\frac ud-\sqrt{u^2+d^2}\right]_0^H\\
&=2\left[H\,\operatorname{arsinh}\frac Hd-\sqrt{H^2+d^2}+d\right].
\end{aligned}\]
当$H\gg a$时，
\[K_H(d)=2H\left(\ln\frac{2H}{d}-1\right)+2d+O(d^2/H).
\]
对数项的角积分可分部积分：
\[\begin{aligned}
\int_0^{2\pi}\cos\phi\ln[2\sin(\phi/2)]\dd\phi
&=\left[\sin\phi\ln[2\sin(\phi/2)]\right]_0^{2\pi}
-\frac12\int_0^{2\pi}\sin\phi\cot\frac\phi2\dd\phi\\
&=-\int_0^{2\pi}\cos^2(\phi/2)\dd\phi
=-\left[\frac\phi2+\frac{\sin\phi}{2}\right]_0^{2\pi}=-\pi.
\end{aligned}\]
分部积分的端点项按极限取值，因$\phi\ln\phi\to0$而消失；中间用了$\sin\phi=2\sin(\phi/2)\cos(\phi/2)$。
后一角积分令$v=\phi/2$，再令$w=\cos v$，可得
\[\begin{aligned}
\int_0^{2\pi}\cos\phi\sin(\phi/2)\dd\phi
&=2\int_0^\pi(1-2\sin^2v)\sin v\dd v\\
&=2\int_{-1}^1(2w^2-1)\dd w
=2[2w^3/3-w]_{-1}^1=-\frac43.
\end{aligned}\]
代入两项，得到
\[\boxed{L=\mu_0n^2\left[\pi a^2H-\frac83a^3
+O(a^4/H)\right]}.
\]
图~\ref{fig:i2-e1-s}右图的电感比值为$L/L_\infty=1-8a/(3\pi H)+O(a^2/H^2)$，首阶端部修正为负。$H$与$a$可比时须计算原积分，不能继续使用长管展开。这里采用密绕电流片的几何电感模型，离散线径、匝间距及导体内部磁能另计。
\end{solution}

\begin{example}[为什么无限单根导线没有有限自感]\label{ex:i2-e2}
半径$a$的无限长直导线通电流$I$，没有给出返回电流。暂在外半径$b$处截断场能，求长度$H$内的外部磁能和对应电感，再讨论$b\to\infty$。 见图~\ref{fig:i2-e2-q}。
\end{example}

\begin{figure}[H]\centering
\begin{tikzpicture}[>=Stealth,line cap=round,line join=round]
\draw[line width=1pt](0,-1.6)--(0,1.6)node[above]{$I$};\draw[->](0,-1.1)--(0,-.2);
\draw[axis](0,0)--(3,0)node[right]{$r$};\draw[guide](0,0)ellipse(.5 and .16);\draw[guide](0,0)ellipse(2.3 and .6);
\node[below]at(.5,0){$a$};\node[below]at(2.3,0){$b$};\node[right]at(1.2,1.25){观察长度$H$};
\end{tikzpicture}
\caption{外半径b暂作为积分边界，尚未代表真实返回电流。}\label{fig:i2-e2-q}
\end{figure}
\begin{solution}
图~\ref{fig:i2-e2-q}中，导线外的磁场为$B_\theta=\mu_0I/(2\pi r)$。只计算半径$a<r<b$的磁能，长度为$H$的体积元为$2\pi rH\dd r$，所以
\[W_{\rm out}=\int_a^b\frac{B_\theta^2}{2\mu_0}\,2\pi rH\dd r
=\frac{\mu_0I^2H}{4\pi}\int_a^b\frac{\dd r}r
=\frac{\mu_0I^2H}{4\pi}[\ln r]_a^b
=\frac{\mu_0I^2H}{4\pi}\ln\frac ba.
\]
相应的外部电感为$L_{\rm out}/H=\mu_0\ln(b/a)/(2\pi)$。图~\ref{fig:i2-e2-s}的对数曲线没有有限的远区极限。


\begin{figure}[H]\centering
\begin{tikzpicture}[>=Stealth,line cap=round,line join=round]
\draw[axis](-.1,0)--(4.5,0)node[right]{$b/a$};\draw[axis](0,-.1)--(0,2.5)node[above]{$L_{\rm out}/H$};
\draw[curve,domain=1:4.1,samples=71]plot(\x,{ln(\x)});
\node[right]at(2.2,.45){$\propto\ln(b/a)$};
\end{tikzpicture}
\caption{远区贡献按对数增长，单根导线的电感不能脱离闭合路径定义。}\label{fig:i2-e2-s}
\end{figure}

令$b\to\infty$时，外磁能和外电感都发散。有限半径$a$只处理了近导线的奇点，不能消除远处的对数发散。若电流均匀分布在实心截面，导线内部还增加有限的$\mu_0H/(8\pi)$，也不改变这一结论。

因此“单根无限导线的自感”必须同时给出回流路径或远区截断。实际供电回路包含返回电流，其远场与去程电流的场相抵消，完整回路因而具有有限磁能。也不能把选择磁矢势零点的参考长度直接解释成一个真实回流半径；前者是势的表示约定，后者改变物理磁场与能量。
\end{solution}

\begin{example}[有限截面双线回路的完整磁链]\label{ex:i2-e3}
两根平行圆柱导线半径均为$a$，轴距$b>2a$，长$H\gg a,b$，分别通大小相等、方向相反的电流$I$。电流在各自截面均匀分布，导体与周围均取$\mu_0$。求单位长度电感，保留导线内部磁链，并说明为什么表面间的一条切面磁通不能直接给出全部答案。 见图~\ref{fig:i2-e3-q}。
\end{example}

\begin{figure}[H]\centering
\begin{tikzpicture}[>=Stealth,line cap=round,line join=round]
\draw[curve](-1.35,0)circle(.55);\draw[curve](1.35,0)circle(.55);\node at(-1.35,0){$\odot I$};\node at(1.35,0){$\otimes I$};
\draw[<->](-1.35,-1)--(1.35,-1)node[midway,below]{$b$};\draw[->](-1.35,0)--(-1.35,.55)node[midway,left]{$a$};
\node[above]at(0,1.2){两导线内电流密度均匀};\draw[axis](-2.2,-1.7)--(2.2,-1.7)node[right]{$x$};
\end{tikzpicture}
\caption{端部回接不影响足够长中段的单位长度结果。}\label{fig:i2-e3-q}
\end{figure}
\begin{solution}
图~\ref{fig:i2-e3-q}采用轴距$b$，不是两表面的间隙；条件为$b>2a$，并暂忽略两端。正、负两导线中$j_z$分别为$\pm I/(\pi a^2)$。下面用磁矢势的截面平均计算完整磁链，避免任取一条穿线截面的路径。

对正电流圆柱，外部$A_z=-\mu_0I\ln(\rho/\rho_0)/(2\pi)$。内部满足$-\partial_\rho A_z=\mu_0I\rho/(2\pi a^2)$，由表面$a$积分到$\rho$，
\[A_z(\rho)-A_z(a)=-\frac{\mu_0I}{2\pi a^2}\int_a^\rho s\dd s
=-\frac{\mu_0I}{2\pi a^2}\left[\frac{s^2}{2}\right]_a^\rho
=\frac{\mu_0I}{4\pi}\left(1-\frac{\rho^2}{a^2}\right).\]
用表面连续性代入$A_z(a)$，得到
\[A_{\rm self}(\rho)=-\frac{\mu_0I}{2\pi}\ln\frac a{\rho_0}
+\frac{\mu_0I}{4\pi}\left(1-\frac{\rho^2}{a^2}\right).
\]
截面平均为
\[\langle\rho^2\rangle=\frac1{\pi a^2}\int_0^a\rho^2\,2\pi\rho\dd\rho
=\frac2{a^2}[\rho^4/4]_0^a=\frac{a^2}2.\]
所以$\langle A_{\rm self}\rangle=-\mu_0I\ln(a/\rho_0)/(2\pi)+\mu_0I/(8\pi)$。图~\ref{fig:i2-e3-s}右图的抛物线正是内部矢势。


\begin{figure}[H]\centering
\begin{tikzpicture}[>=Stealth,line cap=round,line join=round]
\begin{scope}[xshift=-2.5cm]\draw[curve](0,0)circle(1);\draw[guide](0,0)circle(.65);\draw[->](0,0)--(.65,0)node[above]{$\rho$};\node at(-.4,.45){$\odot$};
\node[below]at(0,-1.4){内部磁通只链接部分电流};\end{scope}
\begin{scope}[xshift=2.7cm]\draw[axis](0,0)--(2.3,0)node[right]{$\rho/a$};\draw[axis](0,0)--(0,1.7)node[above]{$A_z$};
\draw[curve,domain=0:1,samples=41]plot(\x,{1.4-.6*(\x)^2});\draw[guide](1,0)--(1,.8);\node[below]at(1,0){$1$};\node[right]at(1.1,1.35){截面平均};\end{scope}
\end{tikzpicture}
\caption{圆柱内部自场需作截面平均，邻导线的调和场可用平均值性质。}\label{fig:i2-e3-s}
\end{figure}

另一根负电流导线在正导线截面内产生调和的$A_z$，因为其源在圆盘外。调和函数的圆盘平均等于圆心值，所以
\[\langle A_{\rm other}\rangle=\frac{\mu_0I}{2\pi}\ln\frac b{\rho_0}.
\]
两项相加为$\langle A_+\rangle=\mu_0I\ln(b/a)/(2\pi)+\mu_0I/(8\pi)$；由对称性，负导线的平均值为$-\langle A_+\rangle$。

单位长度磁能为
\[w_B=\frac12\int j_zA_z\dd A
=\frac12[I\langle A_+\rangle-I\langle A_-\rangle]
=\frac{\mu_0I^2}{2\pi}\ln\frac ba+\frac{\mu_0I^2}{8\pi}.
\]
因此完整双线回路的单位长度电感为
\[\boxed{L'=\frac{\mu_0}{\pi}\left(\ln\frac ba+\frac14\right)}.
\]
其中两根导线的内部自感各为$\mu_0/(8\pi)$，合为$\mu_0/(4\pi)$。在这个均匀体电流、圆截面、无限长二维模型中，调和平均性质使上式对所有$b>2a$成立。

若只把两个表面之间的磁通积分到$x=b-a$，会得到含有表观一阶$a/b$修正的$\ln[(b-a)/a]$。图~\ref{fig:i2-e3-s}左图中，导体内部各电流丝的链接权重不同，邻线场也穿入整个截面，这两部分须一起计入。若改成高频表面电流，电流沿圆周重新分布，其外感为$\mu_0\operatorname{arcosh}[b/(2a)]/\pi$；推导见习题~\ref{ex:i2-h4}。该模型没有本题均匀体电流的内部电感常数。
\end{solution}

\begin{example}[同轴表面电流、电容与传播参数]\label{ex:i2-e4}
无限长同轴导体的内、外电流面半径为$a,b$，$b>a$，两面通相反电流$\pm I$，电流仅在表面。间隙介质的参数为$\epsilon,\mu$。求单位长度电感与电容，并说明$L'C'$关系及改成实心体电流时的变化。 见图~\ref{fig:i2-e4-q}。
\end{example}

\begin{figure}[H]\centering
\begin{tikzpicture}[>=Stealth,line cap=round,line join=round]
\fill[even odd rule,gray!10](0,0)circle(1.6)(0,0)circle(.65);\draw[curve](0,0)circle(.65);\draw[curve](0,0)circle(1.6);
\node at(0,0){$\odot I$};\foreach\a in{30,120,210,300}{\node at(\a:1.6){$\otimes$};}
\draw[->](0,0)--(40:.65)node[midway,above]{$a$};\draw[->](0,0)--(-30:1.6)node[midway,below]{$b$};\node[right]at(1.9,.65){$\epsilon,\mu$};
\end{tikzpicture}
\caption{两种场均位于同轴间隙，场的方向彼此垂直。}\label{fig:i2-e4-q}
\end{figure}
\begin{solution}
图~\ref{fig:i2-e4-q}的内、外导体表面分别通$+I,-I$，外导体以外包围的总电流为零。由Ampere定律，磁场只在$a<r<b$的介质间隙中存在：
\[B_\theta(r)=\frac{\mu I}{2\pi r}.
\]
图~\ref{fig:i2-e4-s}左图的磁能积分为
\[w_B=\int_a^b\frac{B_\theta^2}{2\mu}\,2\pi r\dd r
=\frac{\mu I^2}{4\pi}\int_a^b\frac{\dd r}r
=\frac{\mu I^2}{4\pi}[\ln r]_a^b
=\frac{\mu I^2}{4\pi}\ln\frac ba,
\qquad
\boxed{L'=\frac{\mu}{2\pi}\ln\frac ba}.
\]
按链接磁通积分也得同一结果，因为这里电流在表面，整个间隙磁通都链接内导体电流，不存在体电流的部分链接权重。


\begin{figure}[H]\centering
\begin{tikzpicture}[>=Stealth,line cap=round,line join=round]
\begin{scope}[xshift=-2.5cm]\draw[curve](0,0)circle(.65);\draw[curve](0,0)circle(1.5);\draw[->](1,0)arc(0:75:1)node[above]{$\vct B$};\node at(0,0){$+I$};\node[below]at(0,-1.8){磁能在间隙内};\end{scope}
\begin{scope}[xshift=2.5cm]\draw[curve](0,0)circle(.65);\draw[curve](0,0)circle(1.5);\foreach\a in{0,90,180,270}{\draw[->](\a:.7)--(\a:1.4);}\node[below]at(0,-1.8){$\vct E$沿径向};\end{scope}
\end{tikzpicture}
\caption{磁场沿圆周、电场沿半径，两种能量积分分别给出L与C。}\label{fig:i2-e4-s}
\end{figure}

若内导体单位长度带自由电量$+\lambda_e$、外导体带$-\lambda_e$，右图中的电场为$E_r=\lambda_e/(2\pi\epsilon r)$。两端电压
\[V=\int_a^bE_r\dd r=\frac{\lambda_e}{2\pi\epsilon}[\ln r]_a^b
=\frac{\lambda_e}{2\pi\epsilon}\ln\frac ba,
\qquad
\boxed{C'=\frac{2\pi\epsilon}{\ln(b/a)}}.
\]
所以$L'C'=\mu\epsilon$。在无损、表面电流、横电磁波的长线近似下，传播速率为$1/\sqrt{L'C'}=1/\sqrt{\mu\epsilon}$。

若内导体改成均匀体电流的实心圆柱，还须加内部电感$\mu_{\rm in}/(8\pi)$。这时电流分布和导体内部场已不属于理想表面电流的TEM模型，不能一边加入内部电感，一边仍无条件要求$L'C'=\mu\epsilon$。外回流壳若有有限厚度和体电流，也要按其内部链接另作积分。
\end{solution}


\section{互感、互易与耦合强度}\label{sec:i2-mutual}
互感的代数符号取决于两回路的正方向，互易关系来自线性磁场的能量结构。绕向确定后，磁能非负性还给出耦合强度的上界。

\subsection{有向磁链与Neumann公式}
两个固定回路各选定正向电流$I_1,I_2$及相应曲面法向。线性叠加给
\[\lambda_1=L_1I_1+M_{12}I_2,\qquad
\lambda_2=M_{21}I_1+L_2I_2.
\]
磁场线实际怎样弯曲，并不直接决定$M$的代数符号；应看正电流在另一回路上产生的有向磁通。反转一个回路正向后，该回路电流变量和两个互感系数都变号，物理预测不变。

真空中细线回路$C_2$的矢势为
\[\vct A_2(\vct r_1)=\frac{\mu_0I_2}{4\pi}
\oint_{C_2}\frac{\dd\vct l_2}{|\vct r_1-\vct r_2|}.
\]
用Stokes定理求它链接$C_1$的磁通，得到
\[\boxed{M_{12}=\frac{\mu_0}{4\pi}
\oint_{C_1}\oint_{C_2}\frac{\dd\vct l_1\cdot\dd\vct l_2}{|\vct r_1-\vct r_2|}}.
\]
交换两积分次序，分母和点积不变，故$M_{12}=M_{21}=M$。这里两回路不应有未处理的重合细线段；自感或共边网孔需要保留有限线径。线性、互易铁芯中仍有互易关系，但不能继续把真空核直接用于铁芯内部。

\subsection{充磁顺序与交叉能量}
先将$I_1$建立到给定值、再建立$I_2$，第一步输入$L_1I_1^2/2$。第二步中第二源输入$L_2I_2^2/2$；同时第一回路的磁链增加$M\dd I_2$，第一源还输入$I_1M\dd I_2$。总能量为
\[W_B=\frac12L_1I_1^2+\frac12L_2I_2^2+MI_1I_2.
\]
倒过来建立电流时，交叉项变为$M_{21}I_1I_2$。能量与建立顺序无关，给出同一互易条件。不能只计算正在改变电流的那个电源，漏掉维持另一电流的源功。

无源线性系统磁能非负。将上式配方，
\[W_B=\frac{L_1}{2}\left(I_1+\frac{M}{L_1}I_2\right)^2
+\frac12\left(L_2-\frac{M^2}{L_1}\right)I_2^2,
\]
从而$M^2\le L_1L_2$。定义耦合系数$k=M/\sqrt{L_1L_2}$，有$|k|\le1$。等号对应某个电流组合完全不产生磁能的理想极限，并不表示任意几何的两个实际线圈都能达到。


\begin{example}[嵌套长短螺线管的互感]\label{ex:i2-e5}
两个同轴密绕长螺线管半径$a_1>a_2$、长度$H_1\ge H_2$、匝数密度$n_1,n_2$。内管完全置于外管中部，各自长度远大于半径。求自感、互感与耦合系数，讨论半径相同但长度不同的情形。 见图~\ref{fig:i2-e5-q}。
\end{example}

\begin{figure}[H]\centering
\begin{tikzpicture}[>=Stealth,line cap=round,line join=round]
\draw[curve](-2.5,-.9)rectangle(2.5,.9);\draw[curve](-1.45,-.5)rectangle(1.45,.5);
\foreach\x in{-2.2,-1.7,-1.2,-.7,-.2,.3,.8,1.3,1.8,2.3}{\draw(\x,0)ellipse(.14 and .9);}
\draw[->](-.8,0)--(.8,0);\node[above]at(0,1.15){$n_1,H_1,a_1$};\node[below]at(0,-.15){$n_2,H_2,a_2$};
\draw[<->](-1.45,-1.2)--(1.45,-1.2)node[midway,below]{$H_2$};
\end{tikzpicture}
\caption{内管截面和内管长度决定共同链接的场。}\label{fig:i2-e5-q}
\end{figure}
\begin{solution}
图~\ref{fig:i2-e5-q}中内管完全位于外管内部，且两管均足够长，忽略各自端场。取两正电流所生轴向磁场同向。先只让外管通$I_1$，其内场为$B_1=\mu_0n_1I_1$。内管每匝只链接图~\ref{fig:i2-e5-s}中的面积$S_2=\pi a_2^2$，共有$n_2H_2$匝，故
\[\lambda_{2\leftarrow1}=n_2H_2S_2\mu_0n_1I_1,\qquad
\boxed{M=\mu_0n_1n_2S_2H_2}.
\]
不能把外管截面$S_1$或外管全长$H_1$同时代入互感，因为只有两者共同链接的那部分场参与。


\begin{figure}[H]\centering
\begin{tikzpicture}[>=Stealth,line cap=round,line join=round]
\draw[curve](0,0)circle(1.5);\draw[curve](0,0)circle(.85);\fill[gray!15](0,0)circle(.85);
\draw[->](0,0)--(45:.85)node[midway,left]{$a_2$};\draw[->](0,0)--(-25:1.5)node[midway,below]{$a_1$};
\node[right]at(1.9,.6){链接外场的内管面积$S_2$};\node[right]at(1.9,-.1){$k^2=S_2H_2/(S_1H_1)$};
\end{tikzpicture}
\caption{共同磁能占比由几何面积与重叠长度一起决定。}\label{fig:i2-e5-s}
\end{figure}

自感分别为$L_1=\mu_0n_1^2S_1H_1$、$L_2=\mu_0n_2^2S_2H_2$，于是
\[\boxed{k=\sqrt{\frac{S_2H_2}{S_1H_1}}
=\frac{a_2}{a_1}\sqrt{\frac{H_2}{H_1}}}.
\]
图~\ref{fig:i2-e5-s}把面积比与长度比明确分开。即使两管半径相同，若内管较短，仍有外管的一部分自场不链接内管，因此$k<1$。

反向校验时，让内管通$I_2$：仅在内管所占长度$H_2$和面积$S_2$内有近似均匀场，外管在这段长度有$n_1H_2$匝，结果仍为同一个$M$。若反转一个绕向，$M$和$k$都变负，自感不变。若内管部分伸出，首阶模型还应把$H_2$换成实际重叠长度，端部过渡则不再是严格的突变阶跃。
\end{solution}


\section{耦合电感的连接与退化情形}\label{sec:i2-circuit}
串联、并联接线都保留原线圈间的互感。各支路的磁链关系与端口约束共同确定等效电感；完全耦合时，矩阵可能奇异，需要回到原方程判断允许的电压和电流。

用同名端标出互感符号，取两个电流都从同名端流入时$M>0$。在几何固定的情况下，端口电压按无源符号约定为$V_1=L_1\dot I_1+M\dot I_2$、$V_2=M\dot I_1+L_2\dot I_2$。感应电动势则与这些端口压降反号。列电路方程前应选定使用哪一个量，不能把同一项的符号在中途切换。

串联时同一电流通过两绕组，互感项相加或相减。并联时端口电压相同，支路电流变化率需要联立求解；磁耦合会改变分流比例，不能只按$L_1,L_2$各自的倒数分配。初始电流还可能包含不被端口直接改变的环流分量。


\begin{example}[串并联等效电感及完全耦合]\label{ex:i2-e6}
两线圈自感为$L_1,L_2$，互感绝对值为$M$。分别求串联同向、串联反向、同名端并联和反接并联的等效电感。再讨论$M^2=L_1L_2$时哪些公式需要另作处理。 见图~\ref{fig:i2-e6-q}。
\end{example}

\begin{figure}[H]\centering
\begin{tikzpicture}[>=Stealth,line cap=round,line join=round]
\begin{scope}[yshift=.8cm]\draw[curve](-2,0)--(-1.4,0);\draw[decorate,decoration={coil,aspect=.5,segment length=4pt,amplitude=3pt}](-1.4,0)--(-.2,0);
\draw[curve](-.2,0)--(.3,0);\draw[decorate,decoration={coil,aspect=.5,segment length=4pt,amplitude=3pt}](.3,0)--(1.5,0);\draw[curve](1.5,0)--(2,0);
\node[above]at(-.8,0){$L_1$};\node[above]at(.9,0){$L_2$};\node[below]at(0,-.2){串联，互感$\pm M$};\end{scope}
\begin{scope}[yshift=-1cm]\draw[curve](-2,0)--(-1.5,0)--(-1.5,.4)--(1.5,.4)--(1.5,0)--(2,0);
\draw[curve](-1.5,0)--(-1.5,-.4)--(1.5,-.4)--(1.5,0);\node[fill=white]at(0,.4){$L_1$};\node[fill=white]at(0,-.4){$L_2$};\node[below]at(0,-.75){并联，同端电压};\end{scope}
\end{tikzpicture}
\caption{连接方式改变互感符号；并联必须先解支路方程。}\label{fig:i2-e6-q}
\end{figure}
\begin{solution}
图~\ref{fig:i2-e6-q}上图串联同向时$I_1=I_2=I$，端口电压相加为$(L_1+L_2+2M)\dot I$；反接一个绕组时将$M$换成$-M$，所以
\[\boxed{L_{\rm ser}^{\pm}=L_1+L_2\pm2M}.
\]
由$|M|\le\sqrt{L_1L_2}$知两式均非负，极限最小值为$(\sqrt{L_1}-\sqrt{L_2})^2$。

并联同名端时，两支路电压均为$V$，原方程为
\[L_1\dot I_1+M\dot I_2=V,\qquad
M\dot I_1+L_2\dot I_2=V.\]
记$D=L_1L_2-M^2>0$。第一式乘$L_2$减去第二式乘$M$，得$D\dot I_1=(L_2-M)V$；第二式乘$L_1$减去第一式乘$M$，得$D\dot I_2=(L_1-M)V$。于是
\[\dot I_1=\frac{L_2-M}{D}V,\qquad
\dot I_2=\frac{L_1-M}{D}V.
\]
端口总电流$I=I_1+I_2$，因此
\[\boxed{L_{\rm par}=\frac{L_1L_2-M^2}{L_1+L_2-2M}}.
\]
反接并联时将$M$换成$-M$。在$D>0$时，端口电压决定两支路的电流变化率，支路初值还须另给。图~\ref{fig:i2-e6-s}左图取$L_1=L_2=L$，此时$\dot I_1=\dot I_2=V/(L+M)$，故$I_1-I_2$保持初值。端口等效电感没有包含这项环流信息。


\begin{figure}[H]\centering
\begin{tikzpicture}[>=Stealth,line cap=round,line join=round]
\begin{scope}
\draw[axis](0,0)--(3,0)node[right]{$I_1$};\draw[axis](0,0)--(0,2.6)node[above]{$I_2$};
\draw[curve](.25,.8)--(1.9,2.45);
\fill(.6,1.15)circle(1.4pt)node[left]{$P_0$};
\draw[->,thick](.9,1.45)--(1.4,1.95);
\node[align=center]at(1.5,-.65){$M<L$：$I_1-I_2$为常数\\初值确定所在直线};
\end{scope}
\begin{scope}[xshift=5cm]
\draw[axis](0,0)--(3,0)node[right]{$I_1$};\draw[axis](0,0)--(0,2.6)node[above]{$I_2$};
\draw[curve](.15,2.35)--(2.35,.15);
\draw[<->,thick](.85,1.65)--(1.65,.85);
\node[above right]at(1.3,1.3){$I_1+I_2=I(t)$};
\node[align=center]at(1.5,-.65){$M=L$：给定总电流时\\分流仍可沿线改变};
\end{scope}
\end{tikzpicture}
\caption{$L_1=L_2=L$的同名端并联。左图为非退化方程决定的轨迹；右图为完全耦合时某一时刻的分流约束。}\label{fig:i2-e6-s}
\end{figure}

完全耦合时，原电感矩阵奇异，须回到支路方程。若$L_1=L_2=M=L$，两式都化为$V=L(\dot I_1+\dot I_2)=L\dot I$，因此同向并联的等效电感为$L$；原分式在此是$0/0$，应由支路方程确定其值。两个方程完全相同，已没有决定$\dot I_1-\dot I_2$的独立方程。图~\ref{fig:i2-e6-s}右图中，总电流已知后，分流仍可沿约束直线改变，初始分流也不能确定其后的变化规律。求唯一的支路解须保留漏感、绕组电阻等实际因素。若先在$M<L$下求解再取极限，所得分流取决于所保留的因素及其趋零方式。

若$M=\sqrt{L_1L_2}$且$L_1\ne L_2$，两绕组的理想电压比固定为$\sqrt{L_1/L_2}$，而直接并联又要求电压相等。因此非零端口电压与理想模型不相容；公式趋于零反映的是零磁能模和很大的环流响应，实际由漏感、绕组电阻限制。这个退化极限不能解释为任意两个完全耦合线圈都能无条件并联。

完全耦合后若反接并联，令$a=\sqrt{L_1}$、$b=\sqrt{L_2}$，两个支路方程为
\[V=a(a\dot I_1-b\dot I_2),\qquad V=-b(a\dot I_1-b\dot I_2).\]
第一式乘$b$、第二式乘$a$后相加，得到$(a+b)V=0$。由于$a,b>0$，端口只能满足$V=0$，并有$a\dot I_1=b\dot I_2$。反接等效电感的零值表示这一理想端口约束，并不允许外加非零电压后仍得到有限的支路变化率。
\end{solution}

\begin{example}[中心小方框旋转产生的双向感应]\label{ex:i2-e7}
边长$l$的大正方形回路通电流$I_1$；其中心放置边长$a\ll l$的小正方形回路。小框绕过中心、位于大框平面内的一条直径以角速度$\omega$转动，法向夹角为$\theta=\omega t$。求首阶互感及两框因转动产生的电动势，再在两电流均恒定时求磁力矩。 见图~\ref{fig:i2-e7-q}。
\end{example}

\begin{figure}[H]\centering
\begin{tikzpicture}[>=Stealth,line cap=round,line join=round]
\draw[curve](-2,-.9)--(1,-.9)--(2,.9)--(-1,.9)--cycle;
\draw[->](-1.5,-.9)--(-.6,-.9)node[midway,below]{$I_1$};
\node[below]at(.2,-1.45){大正方形边长$\ell$};
\draw[guide](-1.7,0)--(1.8,0);
\draw[curve](-.55,-.25)--(.35,-.7)--(.65,.25)--(-.25,.7)--cycle;
\draw[->](0,0)--(0,1.7)node[above]{$\vct B_1$};
\draw[->](0,0)--(.85,1.45)node[right]{$\vct n_2$};
\draw[->](0,.8)arc(90:60:.8)node[above right]{$\theta=\omega t$};
\node[right]at(1.2,-.45){$a,I_2$};
\end{tikzpicture}
\caption{小框中心固定，只改变法向与大框磁场的夹角。}\label{fig:i2-e7-q}
\end{figure}
\begin{solution}
图~\ref{fig:i2-e7-q}中，小框尺寸远小于大框边长，可把大框在小框处的场取为中心值。令$h=l/2$，沿大框某一边以中点为$s=0$，该边到中心的垂距为$h$。Biot--Savart定律给
\[B_{\rm side}=\frac{\mu_0I_1}{4\pi}
\int_{-h}^{h}\frac{h\dd s}{(h^2+s^2)^{3/2}}.
\]
令$t=s/h$、$\dd s=h\dd t$，两端变为$-1,1$；而$\dd[t/\sqrt{1+t^2}]/\dd t=(1+t^2)^{-3/2}$，故
\[\begin{aligned}
B_{\rm side}&=\frac{\mu_0I_1}{4\pi h}
\int_{-1}^{1}\frac{\dd t}{(1+t^2)^{3/2}}\\
&=\frac{\mu_0I_1}{4\pi h}
\left[\frac{t}{\sqrt{1+t^2}}\right]_{-1}^{1}
=\frac{\sqrt2\mu_0I_1}{2\pi l}.
\end{aligned}\]
四边相加，$B_1=2\sqrt2\mu_0I_1/(\pi l)$。小框面积为$a^2$，故互感为
\[\boxed{M(\theta)=M_0\cos\theta,\qquad
M_0=\frac{2\sqrt2\mu_0a^2}{\pi l}}.
\]
图~\ref{fig:i2-e7-s}给出磁链随角度的正弦余弦关系。
\begin{figure}[H]\centering
\begin{tikzpicture}[>=Stealth,line cap=round,line join=round]
\draw[axis](-.1,0)--(4.5,0)node[right]{$\theta$};\draw[axis](0,-.1)--(0,2.5)node[above]{$M/M_0$};
\draw[curve,domain=0:6.283,samples=101]plot({.6*\x},{cos(57.29578*\x)});\draw[guide,domain=0:6.283,samples=101]plot({.6*\x},{sin(57.29578*\x)});
\node[right]at(3.1,1.15){实线：$\cos\theta$};\node[right]at(3.1,.65){虚线：$\sin\theta$};
\end{tikzpicture}
\caption{互感取余弦，感应电动势含角度导数；图中横轴用弧度。}\label{fig:i2-e7-s}
\end{figure}

若两电流由外源分别保持恒定，链式法则给
\[\dot M=M_0(-\sin\theta)\dot\theta=-M_0\omega\sin\theta.\]
转动产生的感应电动势为
\[\boxed{\mathcal E_1=-I_2\dot M=M_0I_2\omega\sin\theta,\qquad
\mathcal E_2=M_0I_1\omega\sin\theta}.
\]
这里同一互感函数同时描述两个方向的感应；“外框静止”并不意味着它链接的磁通不变。

固定两电流时，磁力矩为$\tau_\theta=I_1I_2\partial_\theta M=-M_0I_1I_2\sin\theta$。其功率$I_1I_2\dot M$与磁能交叉项的变化率相同；两个维持电流的电源合计功率为$2I_1I_2\dot M$。若只让一个框开路测电动势，不应又假设其中有同样的恒定电流来算力矩。

大框中心处场梯度由对称性消失，有限小框尺寸的相对修正一般从$O(a^2/l^2)$开始。题中$l$是完整边长；若用半边长作几何参数，中心场式和互感系数都会改变一个因子二，必须与图中的长度约定一致。
\end{solution}

\begin{example}[重叠螺线管的吸力与电源功]\label{ex:i2-e8}
两个长螺线管截面近似相同为$S$，匝数密度为$n_1,n_2$，从相反方向插入，重叠长度为$x$。两电流$I_1,I_2$由独立电源保持不变，且生场同向。忽略端场，求磁能、使重叠增加的磁力及两个电源在小位移中的功。 见图~\ref{fig:i2-e8-q}。
\end{example}

\begin{figure}[H]\centering
\begin{tikzpicture}[>=Stealth,line cap=round,line join=round]
\draw[curve](-2.5,-.7)rectangle(.7,.7);\draw[curve](-.7,-.5)rectangle(2.5,.5);
\draw[->](-2,.15)--(-1.2,.15)node[above]{$B_1$};\draw[->](1.2,.15)--(2,.15)node[above]{$B_2$};
\draw[<->](-.7,-1)--(.7,-1)node[midway,below]{$x$};\node[above]at(-1.8,.85){$n_1,H_1$};\node[above]at(1.8,.85){$n_2,H_2$};\node[below]at(0,-1.65){共同截面$S$};
\end{tikzpicture}
\caption{重叠长度是机械坐标，截面和匝数密度保持给定。}\label{fig:i2-e8-q}
\end{figure}
\begin{solution}
图~\ref{fig:i2-e8-q}中，取$x$为两线圈的重叠长度，增大$x$表示进一步插入。忽略端场时，共同磁通给
\[M(x)=\mu_0n_1n_2Sx.
\]
也可用图~\ref{fig:i2-e8-s}的三个区域积分场能：
\[W_B=\frac{S}{2\mu_0}
\left[(H_1-x)B_1^2+x(B_1+B_2)^2+(H_2-x)B_2^2\right].
\]
展开平方后，独立自能项不随$x$改变，交叉项为$\mu_0n_1n_2SxI_1I_2=MI_1I_2$。


\begin{figure}[H]\centering
\begin{tikzpicture}[>=Stealth,line cap=round,line join=round]
\draw[curve](-2.5,-.55)rectangle(2.5,.55);\draw[guide](-.8,-.8)--(-.8,.8);\draw[guide](.8,-.8)--(.8,.8);
\fill[gray!15](-.8,-.55)rectangle(.8,.55);\node at(-1.65,0){$B_1$};\node at(0,0){$B_1+B_2$};\node at(1.65,0){$B_2$};
\draw[->](-2.1,-1.3)--(-1.1,-1.3);\draw[->](2.1,-1.3)--(1.1,-1.3);
\node[below]at(0,-1.5){相互吸入，重叠长度$x$增大};\node[above]at(0,1.2){$\dd W_{\rm src}=2I_1I_2\dd M$};
\end{tikzpicture}
\caption{重叠区的磁场相加，源功的一半进入磁能，另一半成为机械功。}\label{fig:i2-e8-s}
\end{figure}

保持$I_1,I_2$不变时，两个电源分别补偿磁链变化$\dd\lambda_1=I_2\dd M$与$\dd\lambda_2=I_1\dd M$，所以两源合计输入$2I_1I_2\dd M$。磁能只增加$I_1I_2\dd M$，余下部分转为机械功：
\[F_x\dd x=I_1I_2\dd M,\qquad
\boxed{F_x=\mu_0n_1n_2SI_1I_2}.
\]
两电流生场同向时力使重叠增加；反向时力使重叠减小。外界若要准静态保持该位置，需施加相反的力。

直接写$F_x=-\partial_xW_B|_I$会把方向算反，因为它遗漏维持电流的源功。若两线圈闭合且无电阻，才应固定各自总磁链，对$\frac12\vct\lambda^{\mathsf T}\mathsf L^{-1}(x)\vct\lambda$求负导数，此时电流将随位置变化。实际插入接近线圈端部时，$M(x)$平滑过渡，恒力结果仅是重叠区长度模型的首阶近似。
\end{solution}


\section{磁路、气隙与磁阻网络}\label{sec:i2-magnetic}
磁路近似用分段磁动势和磁通描述连续磁场。漏磁、边缘场及材料非线性的取舍决定近似范围，分支关系则由实际磁路连接确定。

\subsection{边界条件与磁通路径}
磁场始终满足$\nabla\cdot\vct B=0$。跨介质界面取一个薄柱，得到$B_{1n}=B_{2n}$；无自由面电流时再对小矩形使用Ampere定律，得到$H_{1t}=H_{2t}$。在线性介质中$\vct B=\mu\vct H$，故法向磁通密度连续，而切向$B$可随磁导率改变。

若磁场与界面法向的夹角为$\theta_1,\theta_2$，则
\[\frac{\tan\theta_1}{\tan\theta_2}=\frac{\mu_1}{\mu_2}.
\]
磁导率很大的细长铁芯可把主要磁通约束在芯内，但气隙边缘仍会漏磁、展宽。磁路近似把芯内每个截面的场取为沿中线、近似均匀，不能替代边界附近的全部三维场。

\subsection{由Ampere定律定义磁阻}
对一个没有分支的细长磁通管，磁通$\Phi=BS$沿路相同，而$B$可随截面$S(l)$变化。于是
\[\oint H\dd l=\Phi\oint\frac{\dd l}{\mu(l)S(l)}=NI.
\]
定义磁阻与磁动势
\[\mathcal R_m=\oint\frac{\dd l}{\mu S},\qquad
\mathcal F=NI,
\]
便有$\Phi=\mathcal F/\mathcal R_m$。串联各段的磁阻相加；汇合节点的有向磁通相加为零；并联支路承受相同的磁动势降。若线圈链接全部主磁通，
\[\boxed{L=\frac{N^2}{\mathcal R_m}}.
\]
这个$N^2$一份来自磁动势$NI$，另一份来自磁链$N\Phi$。

磁路中的$\Phi$对应电路中的电流，磁动势降对应电压，磁阻对应电阻。对应的是方程结构，并非把磁能损耗解释成电阻热：在线性无磁滞铁芯中，磁阻本身不产生持续焦耳耗散。含$n_v$个节点、$n_e$条支路的连通网络有$n_v-1$个独立磁通连续方程和$n_e-n_v+1$个独立回路方程。多窗口铁芯应选足独立磁通，不能把各窗口都当作同一个$\Phi$。


\begin{example}[环形铁芯的窄气隙]\label{ex:i2-e9}
平均半径$R$的高磁导率环形铁芯，均匀截面为$S$，开有宽$h$的窄气隙，$h\ll\sqrt S\ll R$。铁芯相对磁导率$\mu_r$为常数，绕$N$匝通电流$I$。求磁通密度、自感及闭合气隙的吸力，保留有限$\mu_r$的影响。 见图~\ref{fig:i2-e9-q}。
\end{example}

\begin{figure}[H]\centering
\begin{tikzpicture}[>=Stealth,line cap=round,line join=round]
\fill[even odd rule,gray!16](0,0)circle(1.6)(0,0)circle(1.05);\draw[curve](0,0)circle(1.6);\draw[curve](0,0)circle(1.05);
\fill[white](1,-.15)rectangle(1.8,.15);\draw[curve](1.05,-.15)--(1.6,-.15);\draw[curve](1.05,.15)--(1.6,.15);
\foreach\y in{-.5,-.25,0,.25,.5}{\draw(-1.85,\y)--(-.85,\y+.1);}\node[left]at(-1.85,0){$N,I$};
\draw[<->](1.9,-.15)--(1.9,.15)node[right]{$h$};\node at(0,0){$R$};\node[below]at(0,-2){截面$S$，铁芯磁导率$\mu_0\mu_r$};
\end{tikzpicture}
\caption{气隙短而窄，先忽略漏磁与截面上的场变化。}\label{fig:i2-e9-q}
\end{figure}
\begin{solution}
图~\ref{fig:i2-e9-q}中，设铁芯本身的平均路径长为$\ell_{\rm Fe}\simeq2\pi R$。改变气隙时保持铁芯长度和截面不变。无漏磁近似下，铁芯与气隙链接同一磁通，图~\ref{fig:i2-e9-s}右图的串联磁阻为
\[\mathcal R_m=\frac{\ell_{\rm Fe}}{\mu_0\mu_rS}+\frac h{\mu_0S}.
\]
因此
\[\boxed{B=\frac{\mu_0NI}{h+\ell_{\rm Fe}/\mu_r}},\qquad
\boxed{L=\frac{\mu_0N^2S}{h+\ell_{\rm Fe}/\mu_r}}.
\]
铁芯中$H_{\rm Fe}=B/(\mu_0\mu_r)$，气隙中$H_g=B/\mu_0$；不能因为铁芯磁导率大，就令气隙与铁芯的$H$也相等。


\begin{figure}[H]\centering
\begin{tikzpicture}[>=Stealth,line cap=round,line join=round]
\begin{scope}[xshift=-2.5cm]\draw[curve](-1.2,-.6)rectangle(-.25,.6);\draw[curve](.25,-.6)rectangle(1.2,.6);\foreach\y in{-.35,0,.35}{\draw[->](-.8,\y)--(.8,\y);}\draw[->](-.25,.9)--(.25,.9);\node[above]at(0,.9){$h$};\draw[->](-.5,-1)--(.1,-1)node[right]{$F$闭合气隙};\end{scope}
\begin{scope}[xshift=2.7cm]\draw[curve](-1.5,0)--(1.5,0);\node[draw,fill=white]at(-.65,0){$\mathcal R_{\rm Fe}$};\node[draw,fill=white]at(.8,0){$h/(\mu_0S)$};\node[below]at(0,-.7){两段磁阻串联};\end{scope}
\end{tikzpicture}
\caption{气隙与铁芯共享磁通，但其磁动势降不同。}\label{fig:i2-e9-s}
\end{figure}

固定电流时，对气隙宽度求磁共能导数，
\[F_h=\frac12I^2\frac{\partial L}{\partial h}
=-\frac{\mu_0N^2SI^2}{2(h+\ell_{\rm Fe}/\mu_r)^2}
=\boxed{-\frac{B^2S}{2\mu_0}}.
\]
负号表示力使气隙减小，如图~\ref{fig:i2-e9-s}左图。两端面相互吸引的力各有这个大小；若两端各移动$\dd h/2$，合计机械功仍为$F\,\dd h$，不是再乘一个二。

当$h\gg\ell_{\rm Fe}/\mu_r$时，气隙控制磁阻；当$h\ll\ell_{\rm Fe}/\mu_r$时，铁芯本身限制磁通，不能继续用$B\propto1/h$外推到无穷大。题设$h\ll\sqrt S$保证边缘展宽较小，截面尺度远小于$R$保证铁芯场可近似均匀。若磁场使铁芯进入非线性区，应改用实际的$B(H)$关系，不能保留同一个常数$\mu_r$。
\end{solution}

\begin{example}[分叉磁路的磁通与电感]\label{ex:i2-e10}
一主磁路磁阻为$\mathcal R_c$，返回部分分成两条支路，磁阻比$\mathcal R_1:\mathcal R_2=1:3$。主路上绕$N$匝通电流$I$，忽略漏磁。求两支路磁通、等效磁阻和主绕组电感，并用分段磁能验证。 见图~\ref{fig:i2-e10-q}。
\end{example}

\begin{figure}[H]\centering
\begin{tikzpicture}[>=Stealth,line cap=round,line join=round]
\draw[curve](0,1.5)--(0,.8);\draw[curve](0,.8)--(-1.7,.8)--(-1.7,-1)--(0,-1);\draw[curve](0,.8)--(1.7,.8)--(1.7,-1)--(0,-1);\draw[curve](0,-1)--(0,-1.7);
\node[draw,fill=white]at(0,1.15){$\mathcal R_c$};\node[draw,fill=white]at(-1.7,-.1){$\mathcal R_1$};\node[draw,fill=white]at(1.7,-.1){$\mathcal R_2$};
\draw[->](-1.7,.4)--(-1.7,0);\draw[->](1.7,.4)--(1.7,0);\node[left]at(-2,.5){$\Phi_1$};\node[right]at(2,.5){$\Phi_2$};\node[right]at(.15,1.8){主磁通$\Phi$};
\end{tikzpicture}
\caption{支路有相同磁动势降，主支路磁通为两支路之和。}\label{fig:i2-e10-q}
\end{figure}
\begin{solution}
图~\ref{fig:i2-e10-q}的分叉处没有磁通源，故$\Phi=\Phi_1+\Phi_2$。两返回支路连接同一对磁节点，磁动势降相同：
\[\mathcal R_1\Phi_1=\mathcal R_2\Phi_2.
\]
代入$\mathcal R_1:\mathcal R_2=1:3$，有$\Phi_1=3\Phi_2$，再与通量守恒式联立，得到
\[\boxed{\Phi_1=\frac34\Phi,\qquad\Phi_2=\frac14\Phi}.
\]
图~\ref{fig:i2-e10-s}的面积条表示磁通分配，不是磁通密度分配；若两支路截面不同，还应分别除以截面才能比较$B_1,B_2$。


\begin{figure}[H]\centering
\begin{tikzpicture}[>=Stealth,line cap=round,line join=round]
\draw[axis](0,0)--(4.3,0)node[right]{磁通};\fill[gray!25](0,.3)rectangle(3,1);\fill[gray!10](3,.3)rectangle(4,1);\draw[curve](0,.3)rectangle(4,1);\draw[curve](3,.3)--(3,1);
\node at(1.5,.65){$\Phi_1=3\Phi/4$};\node at(3.5,.65){$\Phi_2$};\node[above]at(2,1.3){$\mathcal R_1:\mathcal R_2=1:3$};
\end{tikzpicture}
\caption{较小的磁阻承担较大的磁通。}\label{fig:i2-e10-s}
\end{figure}

设主绕组$N$匝放在公共支路，全部主磁通均与它链接。并联部分磁阻为$\mathcal R_p=\mathcal R_1\mathcal R_2/(\mathcal R_1+\mathcal R_2)$，所以
\[\Phi=\frac{NI}{\mathcal R_c+\mathcal R_p},\qquad
\boxed{L=\frac{N^2}{\mathcal R_c+\mathcal R_p}}.
\]
能量也可逐段求和：
\[W_B=\frac12\mathcal R_c\Phi^2+
\frac12\mathcal R_1\Phi_1^2+\frac12\mathcal R_2\Phi_2^2
=\frac12(\mathcal R_c+\mathcal R_p)\Phi^2.
\]
为验证并联等效关系，可先不代入$1:3$，由等磁动势降和通量守恒联立得
\[
\Phi_1=\frac{\mathcal R_2}{\mathcal R_1+\mathcal R_2}\Phi,\qquad
\Phi_2=\frac{\mathcal R_1}{\mathcal R_1+\mathcal R_2}\Phi.
\]
于是两个返回支路的能量系数为
\[
\mathcal R_1\left(\frac{\Phi_1}{\Phi}\right)^2
+\mathcal R_2\left(\frac{\Phi_2}{\Phi}\right)^2
=\frac{\mathcal R_1\mathcal R_2^2+\mathcal R_2\mathcal R_1^2}
{(\mathcal R_1+\mathcal R_2)^2}
=\mathcal R_p.
\]
再用$NI=(\mathcal R_c+\mathcal R_p)\Phi$和$\lambda=N\Phi$，总能量化为$NI\Phi/2=LI^2/2$。主绕组链接的是两返回支路的通量之和；仅取其中一支会同时改变磁链和能量。
\end{solution}


\section{变压器中的共同磁通与分支磁通}\label{sec:i2-transformer}
变压器的绕组可能链接共同磁通，也可能只链接某条铁芯分支。电压、电流关系应由这些实际磁链建立，漏磁则另列为未被全部绕组链接的部分。

\subsection{端电压与感应电动势}
绕组按正磁通法向选绕向，感应电动势为$-N\dot\Phi$。若端电压采用电流流入正端的无源约定、忽略绕组电阻，则维持该电流所需的端电压为$V=N\dot\Phi$。下文统一用这种端电压，避免把电动势的负号重复写入回路方程。

理想变压器有两个层面的近似：各绕组完全链接共同磁通，并且磁阻足够小，使维持给定磁通的净安匝可忽略。第一层给出电压与匝数比，第二层才给出负载电流的反比。磁导率有限时还存在励磁电流，瞬时输入、输出功率之差可以储存在磁场中。


\begin{example}[有限磁导率变压器的功率]\label{ex:i2-e11}
两绕组匝数为$N_1,N_2$，完全链接同一线性铁芯磁通$\Phi$，铁芯磁阻为$\mathcal R_m$。忽略绕组电阻和漏磁，一次受交流电压驱动，二次接负载。求电压、电流关系以及瞬时功率差，说明理想变压器极限。 见图~\ref{fig:i2-e11-q}。
\end{example}

\begin{figure}[H]\centering
\begin{tikzpicture}[>=Stealth,line cap=round,line join=round]
\fill[even odd rule,gray!18](-2,-1.4)rectangle(2,1.4)(-1.35,-.8)rectangle(1.35,.8);
\draw[curve](-2,-1.4)rectangle(2,1.4);\draw[curve](-1.35,-.8)rectangle(1.35,.8);
\foreach\y in{-.6,-.3,0,.3,.6}{\draw(-2.25,\y)--(-1.1,\y+.12);}
\node[left]at(-2.25,0){$N_1$};\draw[->](-.8,1.1)--(.8,1.1)node[midway,above]{$\Phi$};\foreach\y in{-.6,-.3,0,.3,.6}{\draw(1.1,\y)--(2.25,\y+.12);}\node[right]at(2.25,0){$N_2$};\node[below]at(0,-1.8){$I_1$入同名端，$I_2$从同名端流向负载};
\end{tikzpicture}
\caption{同名端与电流方向一起决定净励磁安匝。}\label{fig:i2-e11-q}
\end{figure}
\begin{solution}
图~\ref{fig:i2-e11-q}取主、次电流的方向使净励磁安匝为$N_1I_1-N_2I_2$。两绕组链接同一$\Phi$，因此
\[V_1=N_1\dot\Phi,\qquad V_2=N_2\dot\Phi,\qquad
\frac{V_1}{V_2}=\frac{N_1}{N_2}.
\]
这个比值来自共同磁通，不必先令铁芯磁阻为零。有限磁阻满足
\[N_1I_1-N_2I_2=\mathcal R_m\Phi.
\]
所以一次电流由负载反射分量与励磁分量组成：
\[I_1=\frac{N_2}{N_1}I_2+\frac{\mathcal R_m}{N_1}\Phi.
\]
图~\ref{fig:i2-e11-s}中，前一部分供给负载，后一部分建立磁场能。


\begin{figure}[H]\centering
\begin{tikzpicture}[>=Stealth,line cap=round,line join=round]
\draw[curve](-2,0)--(-1.1,0);\draw[->](-2,0)--(-1.35,0)node[above]{$P_1$};\node[draw,fill=white,minimum width=1.8cm,minimum height=.8cm]at(0,0){$W_B=\mathcal R_m\Phi^2/2$};\draw[->](1.1,0)--(2,0)node[above]{$P_2$};\node[below]at(0,-1){$P_1-P_2=\dot W_B$};
\end{tikzpicture}
\caption{瞬时功率差进入或离开铁芯磁能，周期平均可为零。}\label{fig:i2-e11-s}
\end{figure}

输入、输出功率分别为$P_1=V_1I_1$、$P_2=V_2I_2$。相减得到
\[P_1-P_2=\dot\Phi(N_1I_1-N_2I_2)
=\mathcal R_m\Phi\dot\Phi
=\frac{\dd}{\dd t}\frac{\mathcal R_m\Phi^2}{2}.
\]
固定几何、线性无损铁芯中，右端只表示储能变化。交流周期稳态的周期平均为零，所以平均输入等于平均负载功率，瞬时却未必相等。

若$V_1=\widehat V_1\cos\omega t$，由$\dot\Phi=V_1/N_1$积分，
\[\Phi(t)-\Phi(0)=\frac{\widehat V_1}{N_1}\int_0^t\cos\omega s\dd s
=\frac{\widehat V_1}{N_1}\left[\frac{\sin\omega s}{\omega}\right]_0^t
=\frac{\widehat V_1}{N_1\omega}\sin\omega t.\]
无直流偏置时取$\Phi(0)=0$。励磁电流滞后电压四分之一周期；其对应电感为$L_m=N_1^2/\mathcal R_m$。在$\mathcal R_m\to0$的理想极限，励磁电流可忽略，才有$I_1=(N_2/N_1)I_2$及瞬时$P_1=P_2$。若含铁耗、铜耗或漏磁，还须在此之外另列耗散与漏感，不能把它们算入无损磁阻的储能项。
\end{solution}

\begin{example}[共同磁通的三绕组负载折算]\label{ex:i2-e12}
理想铁芯上有三个绕组，匝数$N_1,N_2,N_3$，全部链接同一磁通。一次由$u(t)=U_0\cos\omega t$经电阻$R_1$驱动，两个二次绕组分别接$R_2,R_3$。忽略励磁电流，求一次所见等效负载与各电流，并核对功率。 见图~\ref{fig:i2-e12-q}。
\end{example}

\begin{figure}[H]\centering
\begin{tikzpicture}[>=Stealth,line cap=round,line join=round]
\fill[even odd rule,gray!18](-2,-1.4)rectangle(2,1.4)(-1.35,-.8)rectangle(1.35,.8);
\draw[curve](-2,-1.4)rectangle(2,1.4);\draw[curve](-1.35,-.8)rectangle(1.35,.8);
\foreach\y in{-.6,-.3,0,.3,.6}{\draw(-2.25,\y)--(-1.1,\y+.12);}
\node[left]at(-2.25,0){$N_1$};\draw[->](-.8,1.1)--(.8,1.1)node[midway,above]{$\Phi$};\foreach\y in{.25,.5,.75}{\draw(1.1,\y)--(2.3,\y+.1);}\foreach\y in{-.75,-.5,-.25}{\draw(1.1,\y)--(2.3,\y+.1);}
\node[right]at(2.3,.5){$N_2,R_2$};\node[right]at(2.3,-.5){$N_3,R_3$};\node[below]at(0,-1.8){一次电源$u(t)$，串联$R_1$};
\end{tikzpicture}
\caption{两个二次负载共享磁通，电压按匝数比确定。}\label{fig:i2-e12-q}
\end{figure}
\begin{solution}
图~\ref{fig:i2-e12-q}中三个绕组链接同一个$\Phi$。二次电流均取流出同名端，磁阻忽略后有
\[N_1I_1=N_2I_2+N_3I_3,\qquad
V_j=N_j\dot\Phi,\qquad I_2=\frac{N_2\dot\Phi}{R_2},\quad
I_3=\frac{N_3\dot\Phi}{R_3}.
\]
记$G_N=N_2^2/R_2+N_3^2/R_3$，则$I_1=G_N\dot\Phi/N_1$，一次绕组本身的电压$V_1=N_1\dot\Phi$。因此
\[\boxed{R_{\rm ref}=\frac{V_1}{I_1}
=\frac{N_1^2}{N_2^2/R_2+N_3^2/R_3}}.
\]
图~\ref{fig:i2-e12-s}中两个折算电阻并联，不能把它们直接相加。


\begin{figure}[H]\centering
\begin{tikzpicture}[>=Stealth,line cap=round,line join=round]
\draw[curve](-2.6,0)--(-1.8,0);\node[draw,fill=white]at(-1.4,0){$R_1$};\draw[curve](-1,0)--(-.5,0)--(-.5,.8)--(2,.8)--(2,0)--(2.5,0);
\draw[curve](-.5,0)--(-.5,-.8)--(2,-.8)--(2,0);\node[draw,fill=white]at(.7,.8){$(N_1/N_2)^2R_2$};\node[draw,fill=white]at(.7,-.8){$(N_1/N_3)^2R_3$};
\node[below]at(0,-1.5){负载折算到一次侧后并联};
\end{tikzpicture}
\caption{共同磁通结构使两个折算负载呈并联关系。}\label{fig:i2-e12-s}
\end{figure}

将一次串联电阻也计入，$u=R_1I_1+V_1$，所以
\[\boxed{I_1(t)=\frac{u(t)}{R_1+R_{\rm ref}}},\qquad
\dot\Phi=\frac{N_1u(t)}{N_1^2+R_1G_N}.
\]
将$\dot\Phi$分别代入二次电流，完整结果为
\[
I_2(t)=\frac{N_1N_2u(t)}{R_2(N_1^2+R_1G_N)},\qquad
I_3(t)=\frac{N_1N_3u(t)}{R_3(N_1^2+R_1G_N)}.
\]
记$K_\Phi=N_1U_0/(N_1^2+R_1G_N)$，对$\dot\Phi=K_\Phi\cos\omega t$从零时刻积分，有
\[
\Phi(t)-\Phi(0)=K_\Phi\int_0^t\cos\omega s\dd s
=K_\Phi\left[\frac{\sin\omega s}{\omega}\right]_0^t
=\frac{K_\Phi}{\omega}\sin\omega t.
\]
无直流偏磁的周期解取$\Phi(0)=0$；若初始磁通不为零，理想模型会保留这项常量偏置。

功率核对为$V_1I_1=\dot\Phi(N_2I_2+N_3I_3)=I_2^2R_2+I_3^2R_3$；加上$I_1^2R_1$才等于电源输出。持续直流会在这个理想式中使磁通不断累积，实际铁芯的线性范围将限制可用时间；本题负载折算用于给定的理想交流模型。
\end{solution}

\begin{example}[分磁通铁芯的串联型折算]\label{ex:i2-e13}
双窗口理想铁芯的中柱绕一次$N_1$匝，左、右返回柱分别绕$N_2,N_3$匝并接电阻$R_2,R_3$。一次链接两分支磁通之和，二次分别只链接各自分支。求一次所见等效电阻，说明它与共同磁通三绕组结构的区别。 见图~\ref{fig:i2-e13-q}。
\end{example}

\begin{figure}[H]\centering
\begin{tikzpicture}[>=Stealth,line cap=round,line join=round]
\fill[gray!15](-2,-1.3)rectangle(2,1.3);\fill[white](-1.5,-.8)rectangle(-.35,.8);\fill[white](.35,-.8)rectangle(1.5,.8);
\draw[curve](-2,-1.3)rectangle(2,1.3);\draw[curve](-1.5,-.8)rectangle(-.35,.8);\draw[curve](.35,-.8)rectangle(1.5,.8);
\foreach\y in{-.45,-.15,.15,.45}{\draw(-.5,\y)--(.5,\y+.12);\draw(-2.2,\y)--(-1.25,\y+.12);\draw(1.25,\y)--(2.2,\y+.12);}
\node[above]at(0,1.55){$N_1,\Phi=\Phi_2+\Phi_3$};\node[left]at(-2.2,0){$N_2,R_2$};\node[right]at(2.2,0){$N_3,R_3$};
\end{tikzpicture}
\caption{两窗口有两个独立磁通，主柱磁通为其和。}\label{fig:i2-e13-q}
\end{figure}
\begin{solution}
图~\ref{fig:i2-e13-q}的主绕组包围中柱磁通$\Phi_2+\Phi_3$，两个负载绕组各链接一个返回磁通。它与三个绕组共享一个磁通的结构不同。理想铁芯的两个独立Ampere回路分别给
\[N_1I_1=N_2I_2,\qquad N_1I_1=N_3I_3.
\]
二次的端电压为$V_2=N_2\dot\Phi_2=R_2I_2$、$V_3=N_3\dot\Phi_3=R_3I_3$，而一次电压为
\[V_1=N_1(\dot\Phi_2+\dot\Phi_3)
=\frac{N_1}{N_2}V_2+\frac{N_1}{N_3}V_3.
\]

\begin{figure}[H]\centering
\begin{tikzpicture}[>=Stealth,line cap=round,line join=round]
\draw[curve](-2.5,0)--(2.5,0);\node[draw,fill=white]at(-1.2,0){$(N_1/N_2)^2R_2$};\node[draw,fill=white]at(1.2,0){$(N_1/N_3)^2R_3$};
\node[above]at(0,1){$V_1=N_1(\dot\Phi_2+\dot\Phi_3)$};\node[below]at(0,-1){独立磁通窗口：折算电压相加};
\end{tikzpicture}
\caption{电压因分支磁通相加，折算电阻也呈相加形式。}\label{fig:i2-e13-s}
\end{figure}

消去两个二次电流，得到图~\ref{fig:i2-e13-s}的串联型折算：
\[\boxed{\frac{V_1}{I_1}
=\left(\frac{N_1}{N_2}\right)^2R_2+
\left(\frac{N_1}{N_3}\right)^2R_3}.
\]
这里“串联型”指一次端口的等效电阻相加，并非说两个二次绕组在导线上实际串联。两者可以各自接独立负载。

输入功率为$V_1I_1$，代入电压分解式和安匝关系，逐项得$V_1I_1=V_2I_2+V_3I_3$。如果误令$V_1/N_1=V_2/N_2=V_3/N_3$，就相当于删除了一个独立磁通，得到的是上一题的共同磁通结构。

有限磁阻时，主柱和返回柱还会储能，两个净安匝不再同时为零。应为$\Phi_2,\Phi_3$各列一条磁路方程，再接各自的电路方程；不能只用“变压器匝数比”跳过铁芯的连接关系。
\end{solution}

\section{非线性磁链与有损副边的反作用}
\subsection{割线电感、微分电感和磁能}
在线性无磁滞元件中，磁链$\lambda=Li$，磁能为$Li^2/2$。若铁芯的磁响应非线性但仍单值、可逆，则应从完整的$\lambda(i,x)$关系出发；$x$表示机械位置。固定$x$时有两个不同电感：
\[
 L_{\rm sec}=\frac{\lambda}{i},\qquad
 L_{\rm diff}=\left(\frac{\partial\lambda}{\partial i}\right)_x.
\]
缓慢建立磁场时，电源输送的电磁功为$i\,\dd\lambda$，所以磁能及共能分别为
\[
 W(\lambda,x)=\int_0^\lambda i(\tilde\lambda,x)\,\dd\tilde\lambda,
 \quad
 W'(i,x)=\int_0^i\lambda(\tilde i,x)\,\dd\tilde i.
\]
两者满足$W+W'=i\lambda$。这里$W'$是共能的记号，不是对某个变量求导。
\ExpansionDiagram{fig:dep-i2-nonlinear}{单值磁化曲线下的共能面积及其互补的磁能面积}{
\draw[->](0,0)--(5.3,0)node[right]{$i$};\draw[->](0,0)--(0,3.5)node[above]{$\lambda$};
\fill[gray!10](0,0)rectangle(4,2.88);
\fill[gray!35,domain=0:4,samples=60](0,0)--plot(\x,{.4*\x+.02*\x^3})--(4,0)--cycle;
\draw[thick,domain=0:4.3,samples=60]plot(\x,{.4*\x+.02*\x^3});
\draw[dashed](4,0)--(4,2.88)--(0,2.88);
\node at(2.7,.6){$W'$};\node at(1.4,2){$W$};}
图\ref{fig:dep-i2-nonlinear}的横轴为$i$、纵轴为$\lambda$，故曲线下的面积为$\int\lambda\,\dd i$，即共能；矩形的另一部分才是磁能。它们只有在线性关系下相等。

包含机械位移时，能量平衡为
\[
 \dd W=i\,\dd\lambda-F\,\dd x,\qquad
 \dd W'=\lambda\,\dd i+F\,\dd x.
\]
因此恒磁链取$F=-(\partial W/\partial x)_\lambda$，恒电流取$F=(\partial W'/\partial x)_i$。这两个公式由同一微分式得到，并不对应两种不同的磁力。

\begin{example}[位置相关非线性电感的力和小信号响应]
某可逆、无磁滞磁元件的磁链关系为
\[
 \lambda(i,x)=L(x)i+\alpha(x)i^3,
\]
在讨论范围内$L>0,\alpha>0$。求磁能、共能和沿$x$方向的力；再求固定位置、偏置电流$I$附近小扰动的端口方程。导线电阻为$R$。
\ExpansionDiagram{fig:dep-i2-cubic-q}{机械位置影响整条磁链曲线，电流偏置附近使用局部斜率}{
\draw[->](0,0)--(5,0)node[right]{$i$};\draw[->](0,0)--(0,3.7)node[above]{$\lambda$};
\draw[thick,domain=0:4.2,samples=70]plot(\x,{.35*\x+.025*\x^3});
\draw[dashed,domain=0:3.9,samples=70]plot(\x,{.5*\x+.025*\x^3});
\node[right]at(4,3.5){$x+\dd x$};\fill(2.8,1.5288)circle(2pt);
\draw(1.8,.5908)--(3.8,2.4668);\draw[dashed](2.8,0)--(2.8,1.5288);\node[below]at(2.8,0){$I$};}
\end{example}
\begin{solution}
按图\ref{fig:dep-i2-cubic-q}的横轴从零电流积分，位置$x$保持不变：
\[
 W'(i,x)=\int_0^i[L(x)s+\alpha(x)s^3]\dd s
 =\left[\frac12L(x)s^2+\frac14\alpha(x)s^4\right]_0^i
 =\frac12L(x)i^2+\frac14\alpha(x)i^4.
\]
再由$W=i\lambda-W'$得到
\[
 W=\frac12L(x)i^2+\frac34\alpha(x)i^4.
\]
此式以$i$为参数表示磁能；做恒磁链求导时，$i$仍随$x$改变。为避免漏掉这项，直接使用恒电流共能公式：
\[
 \boxed{F=\frac12\frac{\dd L}{\dd x}i^2+
 \frac14\frac{\dd\alpha}{\dd x}i^4.}
\]
若改用磁能，恒$\lambda$条件给出
\[
 0=(L+3\alpha i^2)\frac{\dd i}{\dd x}
 +L_xi+\alpha_xi^3.
\]
解出$\dd i/\dd x=-(L_xi+\alpha_xi^3)/(L+3\alpha i^2)$。再对参数式$W(i,x)$求全导数，
\[\begin{aligned}
-\left(\frac{\dd W}{\dd x}\right)_\lambda
&=-\frac12L_xi^2-\frac34\alpha_xi^4
-i(L+3\alpha i^2)\frac{\dd i}{\dd x}\\
&=-\frac12L_xi^2-\frac34\alpha_xi^4
+i(L_xi+\alpha_xi^3)
=\frac12L_xi^2+\frac14\alpha_xi^4.
\end{aligned}\]
这与共能公式相同。用恒$i$的$-\partial_xW$会遗漏电源功，连力的符号也可能错误。
\ExpansionDiagram{fig:dep-i2-cubic-a}{偏置附近的割线斜率与微分斜率不同}{
\draw[->](0,0)--(5,0)node[right]{$i$};\draw[->](0,0)--(0,3.6)node[above]{$\lambda$};
\draw[thick,domain=0:4.2,samples=60]plot(\x,{.35*\x+.025*\x^3});
\draw[dashed](0,0)--(4.2,2.2932);\draw(1.8,.5908)--(3.8,2.4668);
\fill(2.8,1.5288)circle(2pt);\node[left]at(1.8,.65){$L_{\rm sec}$};
\node[right]at(3.6,2.3){$L_{\rm diff}$};}
令$i=I+\delta i$，固定$x$，忽略$(\delta i)^2$及更高阶项，
\[
 \delta\lambda=(L+3\alpha I^2)\delta i,\qquad
 \boxed{\delta v=R\,\delta i+
 (L+3\alpha I^2)\frac{\dd\,\delta i}{\dd t}.}
\]
图\ref{fig:dep-i2-cubic-a}中的割线对应$L_{\rm sec}=L+\alpha I^2$，切线对应$L_{\rm diff}=L+3\alpha I^2$。若端口扰动电压为零，方程成为$L_{\rm diff}\delta\dot i+R\delta i=0$。从初值$\delta i_0$积分，
\[\int_{\delta i_0}^{\delta i(t)}\frac{\dd s}s=-\frac{R}{L_{\rm diff}}\int_0^t\dd t',\qquad
[\ln|s|]_{\delta i_0}^{\delta i(t)}=-\frac{Rt}{L_{\rm diff}}.\]
非零解保持初始符号，取指数得
\[\delta i(t)=\delta i_0e^{-Rt/L_{\rm diff}},\]
所以衰减时间为$L_{\rm diff}/R$；零初扰动给恒零解。若机械位置也扰动，还应加上$(L_xI+\alpha_xI^3)\delta\dot x$产生的电压项。这里的磁化关系单值，不包含磁滞损耗。
\end{solution}

\subsection{短路副边使主边同时呈现电抗和电阻}
两线性线圈的位置固定，电流参考方向取使互感项为$+M$，电感矩阵正定：$M^2<L_1L_2$。副边由自身电感$L_2$和串联电阻$R_2>0$闭合，主边电阻为$R_1$。采用$\ee^{\mathrm i\omega t}$相量约定，有
\[
 V_1=(R_1+\mathrm i\omega L_1)I_1+\mathrm i\omega MI_2,
 \quad
 0=(R_2+\mathrm i\omega L_2)I_2+\mathrm i\omega MI_1.
\]
由第二式先解出$I_2=-\mathrm i\omega MI_1/(R_2+\mathrm i\omega L_2)$，再代入第一式。因$(\mathrm i\omega M)(-\mathrm i\omega M)=\omega^2M^2$，得到
\[
 Z_{\rm in}=R_1+\mathrm i\omega L_1+
 \frac{\omega^2M^2}{R_2+\mathrm i\omega L_2}.
\]
把分母乘以共轭，可分别读出主边增加的实部及减少的电感：
\[
 \Delta R=\frac{\omega^2M^2R_2}{R_2^2+\omega^2L_2^2}>0,\quad
 L_{\rm eff}=L_1-\frac{\omega^2M^2L_2}{R_2^2+\omega^2L_2^2}.
\]
有限电阻使副边电流与主边电流之间存在相位差。主边因此同时感受到附加电阻和电感的变化，单用反向电流的实数关系无法描述这一反作用。

\begin{example}[副边耗散怎样反映到主边]
上列耦合线圈满足$L_1=L_2=L$、$M=kL$，$0<k<1$，主边电阻忽略。求$\omega L=R_2$时的输入阻抗、副边电流与主边电流之比，并证明主边获得的平均有功功率全部耗散在$R_2$上。相量幅值统一取有效值。
\ExpansionDiagram{fig:dep-i2-secondary-q}{闭合副边通过互感向主边反映耗散}{
\draw[thick](0,.6)--(0,0)--(1.2,0)--(1.2,.1);
\draw[thick](0,1.4)--(0,2)--(1.2,2)--(1.2,1.9);
\draw[thick](0,1)circle(.4);\draw[domain=-.25:.25,samples=30]plot(\x,{1+.12*sin(720*\x)});
\draw[thick](3,.1)--(3,0)--(4.2,0)--(4.2,.5);
\draw[thick](4.2,.5)--(4.05,.65)--(4.35,.85)--(4.05,1.05)--(4.35,1.25)--(4.2,1.4)--(4.2,2)--(3,2)--(3,1.9);
\foreach \y in {.1,.46,.82,1.18,1.54}{\draw[thick](1.2,\y)arc(-90:90:.18);\draw[thick](3,\y)arc(270:90:.18);}
\fill(1.2,1.98)circle(1.5pt);\fill(3,1.98)circle(1.5pt);
\draw[->](.2,2.12)--(.9,2.12)node[midway,above]{$I_1$};\draw[->](4,2.12)--(3.3,2.12)node[midway,above]{$I_2$};
\node[left]at(-.4,1){$V_1$};\node[right]at(4.45,1){$R_2$};
\draw[<->](1.6,2.3)--(2.7,2.3)node[midway,above]{$M$};\node[below]at(.6,0){$L$};\node[below]at(3.6,0){$L$};}
\end{example}
\begin{solution}
在图\ref{fig:dep-i2-secondary-q}中，两线圈按互感项为正的约定接线。由副边方程及$\omega L=R_2$，
\[
 \frac{I_2}{I_1}=-\frac{\mathrm i\omega kL}{R_2+\mathrm i\omega L}
 =-\frac{k}{2}(1+\mathrm i).
\]
所以$|I_2|^2=k^2|I_1|^2/2$。主边增加的实部为$k^2R_2/2$，有效电感为$L(1-k^2/2)$：
\[
 \boxed{Z_{\rm in}=\frac{k^2R_2}{2}
 +\mathrm iR_2\left(1-\frac{k^2}{2}\right)}.
\]
\ExpansionDiagram{fig:dep-i2-secondary-a}{以主边电流为实轴，副边电流含反相及正交分量}{
\draw[->](-2.2,0)--(2.3,0)node[right]{实部};\draw[->](0,-2)--(0,1.3)node[above]{虚部};
\draw[->,thick](0,0)--(1.7,0)node[above]{$I_1$};
\draw[->,thick](0,0)--(-1.25,-1.25)node[below left]{$I_2$};
\draw[dashed](-1.25,0)--(-1.25,-1.25)--(0,-1.25);
\node[above]at(-1.25,0){$-kI_1/2$};}
图\ref{fig:dep-i2-secondary-a}取$I_1$为实数参考。由于相量为有效值，主边平均输入功率不再乘$1/2$：
\[
 P_{\rm in}=|I_1|^2\operatorname{Re}Z_{\rm in}
 =\frac{k^2R_2}{2}|I_1|^2
 =R_2|I_2|^2.
\]
互感本身没有耗散；主边的实阻抗反映了副边的电阻热。低频$\omega L_2\ll R_2$时，$\Delta R$按$\omega^2$趋零、$L_{\rm eff}\to L_1$；高频$\omega L_2\gg R_2$时，$L_{\rm eff}\to L_1-M^2/L_2>0$。最后的不等式来自电感矩阵正定，保证这个无源回路不会给出负的有效电感。
\end{solution}

\section{习题与解答}

\begin{exercise}[恒电压与变化电感]\label{ex:i2-h1}
理想无电阻线圈的电感为$L(x)=L_0+\alpha x$，$L_0,\alpha>0$。机械部分被外界以$x=ut$匀速拉动，$u>0$；线圈端接恒电压$V$，初始电流为零。求电流、磁力及功率分配。 见图~\ref{fig:i2-h1-q}。
\end{exercise}

\begin{figure}[H]\centering
\begin{tikzpicture}[>=Stealth,line cap=round,line join=round]
\draw[curve](-2,0)--(-.6,0);\node[draw,fill=white]at(0,0){$L(x)=L_0+\alpha x$};\draw[curve](.9,0)--(2,0);
\draw[->](.4,-.8)--(1.4,-.8)node[right]{$x=ut$};\node[above]at(0,.8){恒端电压$V$，$I(0)=0$};
\end{tikzpicture}
\caption{线圈的几何运动由外界指定，端电压保持不变。}\label{fig:i2-h1-q}
\end{figure}
\begin{solution}
图~\ref{fig:i2-h1-q}给定端电压。无电阻回路满足$\dot\lambda=V$，从初态$\lambda(0)=L_0I(0)=0$积分，
\[\lambda(t)-0=\int_0^tV\dd s=V[s]_0^t=Vt.\]
代入$\lambda=LI$与$L(t)=L_0+\alpha ut$，得到
\[\boxed{I(t)=\frac{Vt}{L_0+\alpha ut}}.
\]
图~\ref{fig:i2-h1-s}的电流渐近于$V/(\alpha u)$；这不等于自感电动势消失，而是电流增长变慢的同时，几何磁链项$I\dot L$承担了几乎全部端电压。


\begin{figure}[H]\centering
\begin{tikzpicture}[>=Stealth,line cap=round,line join=round]
\draw[axis](-.1,0)--(4.5,0)node[right]{$t$};\draw[axis](0,-.1)--(0,2.5)node[above]{$I$};
\draw[curve,domain=0:4,samples=71]plot(\x,{1.7*\x/(1+\x)});\draw[guide](0,1.7)--(4,1.7);
\node[above]at(2.5,1.7){$V/(\alpha u)$};
\end{tikzpicture}
\caption{电流趋于有限值时，几何磁链仍持续变化。}\label{fig:i2-h1-s}
\end{figure}

磁力按瞬时固定电流的共能导数为$F_x=\alpha I^2/2$，方向使电感增加。维持给定位移$x=ut$的外力需承担相反的磁力分量，其机械功率为$-F_xu$。

磁能$W_B=V^2t^2/[2(L_0+\alpha ut)]$。先按商法则求电流导数，再按乘积法则求磁能导数：
\[\begin{aligned}
\dot I&=\frac{V(L_0+\alpha ut)-Vt\alpha u}{(L_0+\alpha ut)^2}
=\frac{VL_0}{(L_0+\alpha ut)^2},\\
\dot W_B&=\frac12\dot L I^2+LI\dot I
=\frac12\alpha uI^2+I(V-I\alpha u).
\end{aligned}\]
其中$L\dot I+I\alpha u=V$可由上式直接回代核对。整理得
\[\dot W_B=VI-\frac12\alpha uI^2,\qquad
VI=\dot W_B+F_xu.
\]
电源功率分别用于增加磁能和向机械系统输出功。若误用$V=L\dot I$，会得到不同的电流，且无法同时满足这条功率关系。给定的线性电感规律需在全部位移范围内成立；实际可动铁芯行程有限，不能把渐近曲线无限延长为真实线圈的长期响应。
\end{solution}

\begin{exercise}[长管近似的误差与伸长力]\label{ex:i2-h2}
空气螺线管总匝数$N$、半径$a$固定，长度$H\gg a$。用端部首阶修正估计：忽略端场的相对误差小于$1\%$需多大$H/a$？若保持电流$I$缓慢拉长线圈，再求含首阶修正的轴向磁力。 见图~\ref{fig:i2-h2-q}。
\end{exercise}

\begin{figure}[H]\centering
\begin{tikzpicture}[>=Stealth,line cap=round,line join=round]
\draw[curve](-2,-.75)--(2,-.75);\draw[curve](-2,.75)--(2,.75);
\draw[curve](-2,0)ellipse(.25 and .75);\draw[curve](2,0)ellipse(.25 and .75);
\foreach\x in{-1.7,-1.2,-.7,-.2,.3,.8,1.3,1.8}{\draw(\x,0)ellipse(.23 and .75);}
\draw[axis](-2.7,0)--(2.7,0)node[right]{$z$};\draw[->](-.8,.25)--(.8,.25)node[midway,above]{$\vct B$};
\draw[<->](-2,-1.15)--(2,-1.15)node[midway,below]{$H$};\node[above]at(0,.9){$n,\ I$};\draw[<->](2.5,0)--(2.5,.75)node[midway,right]{$a$};\node[below]at(0,-1.9){固定总匝数$N$，半径$a$不变};
\end{tikzpicture}
\caption{伸长时匝数密度随长度改变。}\label{fig:i2-h2-q}
\end{figure}
\begin{solution}
图~\ref{fig:i2-h2-q}的总匝数$N$固定，故$n=N/H$。将它代入长管含端部修正的公式，得到
\[L(H)=\frac{\mu_0N^2\pi a^2}{H}
\left(1-\frac{8a}{3\pi H}+O(a^2/H^2)\right).
\]
以无限长近似值作分母，首阶相对误差为$\delta_L\simeq8a/(3\pi H)$。要求小于$1\%$，给出
\[\boxed{\frac Ha>\frac{800}{3\pi}\simeq84.88}.
\]
图~\ref{fig:i2-h2-s}只表示这一误差随长径比减小的趋势，不能当作短线圈的精确曲线。


\begin{figure}[H]\centering
\begin{tikzpicture}[>=Stealth,line cap=round,line join=round]
\begin{scope}[x=.04cm,y=30cm]
\draw[axis](0,0)--(115,0)node[right]{$H/a$};
\draw[axis](0,0)--(0,.065)node[above]{$\delta_L$};
\draw[curve,domain=20:110,samples=91]plot(\x,{8/(3*pi*\x)});
\draw[guide](0,.01)--(84.8826,.01)--(84.8826,0);
\node[left]at(0,.01){$1\%$};\node[below]at(84.8826,0){$84.88$};
\node[right]at(35,.048){$\delta_L\simeq\dfrac{8}{3\pi(H/a)}$};
\end{scope}
\end{tikzpicture}
\caption{误差由长径比控制，图示首阶反比规律。}\label{fig:i2-h2-s}
\end{figure}

若固定电流并缓慢增加$H$，半径与$N$保持不变。先展开各幂次再求导，
\[\begin{aligned}
L(H)&=\mu_0N^2\pi a^2\left[H^{-1}-\frac{8a}{3\pi}H^{-2}+O(a^2H^{-3})\right],\\
L'(H)&=\mu_0N^2\pi a^2\left[-H^{-2}+\frac{16a}{3\pi}H^{-3}+O(a^2H^{-4})\right].
\end{aligned}\]
磁力为
\[F_H=\frac12I^2L'(H)
=-\frac{\mu_0N^2\pi a^2I^2}{2H^2}
\left(1-\frac{16a}{3\pi H}+O(a^2/H^2)\right).
\]
主项为负，表示线圈有缩短的趋势。不能在求导时把$n$也视为常数，否则等于在拉伸过程中不断补入新匝，研究的机械系统已经改变。误差不等式只是首阶估计；若要求有严格误差保证，应回到完整有限长度积分，估计下一阶项，而不是把截断符号直接当作精确上界。
\end{solution}

\begin{exercise}[非均匀体电流的内部电感]\label{ex:i2-h3}
半径$a$的长圆柱导体中，轴向电流密度满足$j_z(r)=Ar^p$，$p>-2$，总电流为$I$，磁导率为$\mu$。求单位长度内部电感，并讨论$p=0$与$p\to\infty$两个极限。 见图~\ref{fig:i2-h3-q}。
\end{exercise}

\begin{figure}[H]\centering
\begin{tikzpicture}[>=Stealth,line cap=round,line join=round]
\fill[gray!12](0,0)circle(1.3);\draw[curve](0,0)circle(1.3);\draw[guide](0,0)circle(.7);\draw[->](0,0)--(.7,0)node[above]{$r$};\draw[->](0,0)--(120:1.3)node[midway,left]{$a$};\node[right]at(1.6,.5){$j_z\propto r^p$};\node[right]at(1.6,-.2){$p>-2$};
\end{tikzpicture}
\caption{电流分布作为模型条件给定，外部回流另行处理。}\label{fig:i2-h3-q}
\end{figure}
\begin{solution}
图~\ref{fig:i2-h3-q}的电流密度虽然不均匀，但仍轴对称。总电流由截面积分求出：
\[I=2\pi A\int_0^ar^{p+1}\dd r
=2\pi A\left[\frac{r^{p+2}}{p+2}\right]_0^a
=\frac{2\pi Aa^{p+2}}{p+2}.\]
同样将上限换为$r$即得$I_{\rm enc}=2\pi Ar^{p+2}/(p+2)$，所以
\[\frac{I_{\rm enc}(r)}I=\left(\frac ra\right)^{p+2}.
\]
图~\ref{fig:i2-h3-s}比较了两个指数。靠近表面集中的电流使内部被链接的电流分数减小。


\begin{figure}[H]\centering
\begin{tikzpicture}[>=Stealth,line cap=round,line join=round]
\draw[axis](-.1,0)--(4.5,0)node[right]{$r/a$};\draw[axis](0,-.1)--(0,2.5)node[above]{$I_{\rm enc}/I$};
\draw[curve,domain=0:1,samples=51]plot({3*\x},{1.8*(\x)^2});\draw[guide,domain=0:1,samples=51]plot({3*\x},{1.8*(\x)^4});
\draw[guide](3,0)--(3,1.8);\node[below]at(3,0){$1$};
\draw(-.06,1.8)--(.06,1.8);\node[left]at(0,1.8){$1$};
\node[right,fill=white,inner sep=1pt]at(1.8,.9){$p=0$};
\node[right,fill=white,inner sep=1pt]at(2,.3){$p=2$};
\end{tikzpicture}
\caption{较大的指数把电流移向表面，减少内部磁能。}\label{fig:i2-h3-s}
\end{figure}

Ampere定律给$B_\theta=\mu I_{\rm enc}/(2\pi r)$。单位长度内部磁能为
\[\begin{aligned}
w_{\rm in}&=\int_0^a\frac{B_\theta^2}{2\mu}\,2\pi r\dd r\\
&=\frac{\mu I^2}{4\pi}\int_0^a
\left(\frac ra\right)^{2p+4}\frac{\dd r}{r}
=\frac{\mu I^2}{4\pi a^{2p+4}}\left[\frac{r^{2p+4}}{2p+4}\right]_0^a
=\frac{\mu I^2}{8\pi(p+2)}.
\end{aligned}\]
因此
\[\boxed{L'_{\rm in}=\frac{\mu}{4\pi(p+2)}}.
\]
$p=0$恢复均匀体电流的$\mu/(8\pi)$；$p\to\infty$趋于表面电流，内部电感趋零。也可按磁链权重计算$I^{-1}\int_0^aB_\theta(I_{\rm enc}/I)\dd r$，得到同一式。若把权重遗漏，算的是某一条特殊电流丝的链接，而不是整个分布的电感。题中电流分布可由径向变化的电导率等方式维持，并非假定均匀材料在任意频率下自动具有这个分布。
\end{solution}

\begin{exercise}[双线表面电流的双极坐标解]\label{ex:i2-h4}
两平行理想圆柱导体半径$a$、轴距$b>2a$，间隙为$\epsilon,\mu$，无内部场。已知双极坐标面$\xi=\pm\xi_0$对应这两圆柱，$\cosh\xi_0=b/(2a)$，两尺度因子相同。求单位长度电容，再由TEM关系求电感及远距展开。 见图~\ref{fig:i2-h4-q}。
\end{exercise}

\begin{figure}[H]\centering
\begin{tikzpicture}[>=Stealth,line cap=round,line join=round]
\draw[curve](-1.45,0)circle(.65);\draw[curve](1.45,0)circle(.65);\node at(-1.45,0){$+\lambda_e$};\node at(1.45,0){$-\lambda_e$};
\draw[<->](-1.45,-1.1)--(1.45,-1.1)node[midway,below]{$b$};\node[above]at(0,1.2){表面电流、理想等势圆柱};
\end{tikzpicture}
\caption{相邻表面上的电荷与电流密度都不必均匀。}\label{fig:i2-h4-q}
\end{figure}
\begin{solution}
图~\ref{fig:i2-h4-q}的两圆柱是理想等势面，表面电流会受相邻导线影响而重新分布，不能继续采用均匀体电流模型。用双极坐标，两个圆柱面为$\xi=\pm\xi_0$，圆心距与半径满足$\cosh\xi_0=b/(2a)$，如图~\ref{fig:i2-h4-s}。


\begin{figure}[H]\centering
\begin{tikzpicture}[>=Stealth,line cap=round,line join=round]
\draw[axis](-.1,0)--(4.5,0)node[right]{$b/(2a)$};\draw[axis](0,-.1)--(0,2.5)node[above]{$\xi_0$};
\draw[curve,domain=1.01:4.1,samples=81]plot(\x,{ln(\x+sqrt((\x)^2-1))});
\node[right]at(1.8,.6){$\cosh\xi_0=b/(2a)$};
\end{tikzpicture}
\caption{双极坐标参数随两圆柱距离变化。}\label{fig:i2-h4-s}
\end{figure}

两面电压差为$V$时，取两面电势为$\pm V/2$。因$h_\xi=h_\eta=h$，二维Laplace方程为$\partial_\xi^2\varphi+\partial_\eta^2\varphi=0$。取不含角变量的解，径向方程$\partial_\xi^2\varphi=0$两次积分给$\varphi=A\xi+C$；由$A\xi_0+C=V/2$、$-A\xi_0+C=-V/2$得$A=V/(2\xi_0)$、$C=0$。因此$\varphi=V\xi/(2\xi_0)$；这个满足全部边界条件的解由静电唯一性定理确定。积分法向位移时，两尺度因子相消：
\[\lambda_e=\epsilon\int_0^{2\pi}
\frac{V}{2\xi_0h_\xi}\,h_\eta\dd\eta
=\frac{\epsilon V}{2\xi_0}[\eta]_0^{2\pi}
=\frac{\pi\epsilon V}{\xi_0},\qquad
C'=\frac{\pi\epsilon}{\xi_0}.
\]
无损TEM场的电磁对偶给$L'C'=\mu\epsilon$，因此
\[\boxed{L'=\frac{\mu}{\pi}\operatorname{arcosh}\frac b{2a}}.
\]
远距展开可从反双曲函数的对数式求出。令$\delta=a/b$，则
\[
\operatorname{arcosh}\frac b{2a}
=\ln\frac{b+\sqrt{b^2-4a^2}}{2a}
=\ln\frac ba+\ln\frac{1+\sqrt{1-4\delta^2}}2.
\]
将$\sqrt{1-4\delta^2}=1-2\delta^2-2\delta^4+O(\delta^6)$代入，再用$\ln(1+x)=x-x^2/2+\cdots$，第二个对数为$-\delta^2-3\delta^4/2+O(\delta^6)$。

远离时$\operatorname{arcosh}[b/(2a)]=\ln(b/a)-a^2/b^2+O(a^4/b^4)$，首个几何修正为二次而非一次。

本题导体内部无磁场，因而没有均匀体电流模型中的$\mu/(4\pi)$内部电感。两导线靠近时，两种模型的电流分布与边界条件不同，电感数值也随之不同。上述推导也要求两表面没有接触，$b>2a$，不能跨越几何接触点继续使用同一电容、电感式。
\end{solution}

\begin{exercise}[三线圈的互感是否物理允许]\label{ex:i2-h5}
三个线圈自感均为$L_0>0$，互感为$M_{12}=M_{23}=kL_0$、$M_{13}=0$。求被动线性系统允许的$k$范围，并说明只核对每一对线圈的互感上界为何不够。 见图~\ref{fig:i2-h5-q}。
\end{exercise}

\begin{figure}[H]\centering
\begin{tikzpicture}[>=Stealth,line cap=round,line join=round]
\foreach\x/\t in{-2/1,0/2,2/3}{\node[draw,fill=white,minimum width=1cm,minimum height=.6cm]at(\x,0){$L_0$};\node[above]at(\x,.45){$\t$};}
\draw[<->](-1.4,0)--(-.6,0)node[midway,below]{$kL_0$};
\draw[<->](.6,0)--(1.4,0)node[midway,below]{$kL_0$};
\draw[guide](-2,.9)--(2,.9)node[midway,above]{$M_{13}=0$};
\node[below]at(0,-.95){耦合矩阵的示意连接};
\end{tikzpicture}
\caption{磁场能由全部电流共同确定，邻近耦合不代表可任意指定矩阵。}\label{fig:i2-h5-q}
\end{figure}
\begin{solution}
图~\ref{fig:i2-h5-q}的电感矩阵为
\[\mathsf L=L_0
\begin{pmatrix}1&k&0\\k&1&k\\0&k&1\end{pmatrix}.
\]
各对线圈分别满足$|M|\le L_0$，只给$|k|\le1$；这还不足以保证三个电流同时存在时磁能非负。


\begin{figure}[H]\centering
\begin{tikzpicture}[>=Stealth,line cap=round,line join=round]
\draw[axis](-.1,0)--(4.5,0)node[right]{$k$};\draw[axis](0,-.1)--(0,2.5)node[above]{$\lambda_{\min}/L_0$};
\draw[curve](0,1.7)--(1.7,0)--(2.7,-1);\draw[guide](1.7,-1)--(1.7,1.7);
\node[below right]at(1.7,0){$1/\sqrt2$};\node[right]at(2.2,-.65){负磁能区};
\end{tikzpicture}
\caption{最小特征值决定磁能是否可以对所有电流非负。}\label{fig:i2-h5-s}
\end{figure}

令特征值为$\lambda=L_0x$，将无量纲行列式沿第一行展开，
\[\begin{aligned}
\det(\mathsf L/L_0-x\mathsf 1)
&=(1-x)[(1-x)^2-k^2]-k[k(1-x)]\\
&=(1-x)[(1-x)^2-2k^2].
\end{aligned}\]
令它为零，得$x=1$或$1-x=\pm\sqrt2k$。三个特征值为$L_0$、$L_0(1+\sqrt2k)$、$L_0(1-\sqrt2k)$。因此图~\ref{fig:i2-h5-s}的最小值非负要求
\[\boxed{|k|\le\frac1{\sqrt2}}.
\]
例如$k>1/\sqrt2$时取电流比例$(I_1,I_2,I_3)=I(1,-\sqrt2,1)$，直接代入得到$W_B=2L_0I^2(1-\sqrt2k)<0$，不可能是被动线性磁场的能量。

互感的允许范围由完整磁能矩阵决定；逐对满足上界，仍可能存在使总磁能为负的电流组合。等号时出现一个零磁能模；若它与端口的强迫条件冲突，理想电路会退化，实际漏感或电阻必须参与。矩阵中的零互感是给定的理想几何条件，不应因相邻线圈都有耦合就擅自补入一个未给定的$M_{13}$。
\end{solution}

\begin{exercise}[并联后完全不增流的支路]\label{ex:i2-h6}
两线圈自感为$4L_0,L_0$，互感$M=L_0$，同名端并联并接电压$V(t)$。初始两支路电流均为零。求各支路电流和等效电感，并解释为何一支路可以始终不流电流。 见图~\ref{fig:i2-h6-q}。
\end{exercise}

\begin{figure}[H]\centering
\begin{tikzpicture}[>=Stealth,line cap=round,line join=round]
\draw[curve](-2,0)--(-1.5,0)--(-1.5,1)--(1.5,1)--(1.5,0)--(2,0);\draw[curve](-1.5,0)--(-1.5,-1)--(1.5,-1)--(1.5,0);
\node[draw,fill=white]at(0,1){$4L_0$};\node[draw,fill=white]at(0,-1){$L_0$};\node at(0,0){$M=L_0$};\node[left]at(-2,0){$V(t)$};
\end{tikzpicture}
\caption{两个端电压相同，但支路电流不必按普通倒电感比变化。}\label{fig:i2-h6-q}
\end{figure}
\begin{solution}
图~\ref{fig:i2-h6-q}两支路同端并联，方程为$V=4L_0\dot I_1+L_0\dot I_2$与$V=L_0\dot I_1+L_0\dot I_2$。两式相减，$3L_0\dot I_1=0$，所以初始$I_1=0$时该支路始终没有电流。

第二式变成$V=L_0\dot I_2$，因此
\[\boxed{I_1(t)=0,\qquad I_2(t)=\frac1{L_0}\int_0^tV(s)\dd s,\qquad L_{\rm eq}=L_0}.
\]
图~\ref{fig:i2-h6-s}以电压的时间积分为横坐标；第一支路电流保持为零，端口电流全部流经第二支路。


\begin{figure}[H]\centering
\begin{tikzpicture}[>=Stealth,line cap=round,line join=round]
\draw[axis](0,0)--(3.7,0)node[right]{$Q_V=\int_0^tV\dd t$};\draw[axis](0,0)--(0,2.2)node[above]{$I$};\draw[curve](0,0)--(3,1.8)node[right]{$I_2$};\draw[line width=.7pt](0,0)--(3.4,0)node[below]{$I_1=0$};
\end{tikzpicture}
\caption{互感电压可在零支路电流条件下承担全部端电压。}\label{fig:i2-h6-s}
\end{figure}

没有电流的第一支路仍有磁链$\lambda_1=MI_2$，其变化率$M\dot I_2=V$正好给出端电压。不能由$I_1=0$推断该线圈“没有任何感应”。能量只剩$L_0I_2^2/2$，端口输入$VI_2$恰为其导数，与等效电感一致。

本题$M^2/(L_1L_2)=1/4$，矩阵行列式为$3L_0^2>0$，支路解唯一。若第一支路初始已有电流，它将保持该常量，并改变磁能交叉项；题给零初态才使$I_1$始终为零。
\end{solution}

\begin{exercise}[反向串联的误差敏感性]\label{ex:i2-h7}
两线圈$L_1=L_0(1+\epsilon)^2$、$L_2=L_0$，正互感为$M=(1-\delta)L_0(1+\epsilon)$，其中$|\epsilon|\ll1$、$0\le\delta\ll1$。将它们反向串联，求剩余电感并比较失配与漏耦合的影响。 见图~\ref{fig:i2-h7-q}。
\end{exercise}

\begin{figure}[H]\centering
\begin{tikzpicture}[>=Stealth,line cap=round,line join=round]
\draw[curve](-2,0)--(-1.4,0);\node[draw,fill=white]at(-.7,0){$L_0(1+\epsilon)^2$};\draw[curve](0,0)--(.5,0);\node[draw,fill=white]at(1.2,0){$L_0$};\draw[curve](1.8,0)--(2.3,0);
\node[above]at(0,1){反向串联};\node[below]at(0,-.8){$M=(1-\delta)L_0(1+\epsilon)$};
\end{tikzpicture}
\caption{同样接近理想抵消，两种小偏差进入答案的阶次不同。}\label{fig:i2-h7-q}
\end{figure}
\begin{solution}
图~\ref{fig:i2-h7-q}采用反向串联，所以端口电感为$L_1+L_2-2M$。代入给定参数，不先截断$\epsilon$，得到
\[\begin{aligned}
L_{\rm eq}/L_0
&=(1+\epsilon)^2+1-2(1-\delta)(1+\epsilon)\\
&=\boxed{\epsilon^2+2\delta(1+\epsilon)}.
\end{aligned}\]

\begin{figure}[H]\centering
\begin{tikzpicture}[>=Stealth,line cap=round,line join=round]
\draw[axis](-1.6,0)--(1.8,0)node[right]{$\epsilon$};\draw[axis](0,0)--(0,2)node[above]{$\delta$};
\draw[curve,domain=-1.3:1.3,samples=61]plot(\x,{.8-(\x)^2/3});\node[right]at(.7,1.5){$\epsilon^2+2\delta\simeq$常数};\node[below]at(0,-.5){低剩余电感需同时控制两种偏差};
\end{tikzpicture}
\caption{等值线只表示小参数首项，不替代完整的非负性条件。}\label{fig:i2-h7-s}
\end{figure}

当$|\epsilon|,\delta\ll1$时，首要两项是$\epsilon^2+2\delta$。图~\ref{fig:i2-h7-s}的近似等值线表示：匝数或尺度不匹配造成的影响从$\epsilon^2$开始，而漏耦合造成的影响从$\delta$一阶开始。

若目标$L_{\rm eq}/L_0<\eta\ll1$，就要大致同时满足$|\epsilon|\lesssim\sqrt\eta$与$\delta\lesssim\eta/2$，不能只把两个绕组匝数做得接近。完全耦合$\delta=0$时剩余值为$L_0\epsilon^2$；完全匹配$\epsilon=0$时则为$2\delta L_0$。

在题设的小参数范围内，$\delta\ge0$使耦合系数$k=1-\delta$不超过1；$\epsilon>-1$则保证$\sqrt{L_1}=\sqrt{L_0}(1+\epsilon)$与所给正互感采用同一取支。$L_1$本身是平方，不能把这一条件说成排除负自感。若改用同向串联的加号，互感项增加磁能，便没有这里的低电感抵消。
\end{solution}

\begin{exercise}[转动小线圈的负载阻力矩]\label{ex:i2-h8}
固定一次线圈通恒电流$I_1$，与匀速转动的小线圈互感$M=M_0\cos\omega t$。小线圈电阻$R_2$、自感$L_2$，求交流稳态电流、平均磁阻力矩和平均机械功率。 见图~\ref{fig:i2-h8-q}。
\end{exercise}

\begin{figure}[H]\centering
\begin{tikzpicture}[>=Stealth,line cap=round,line join=round]
\draw[axis](-2.6,0)--(2.9,0)node[right]{$z$};
\draw[curve](-1.3,0)ellipse(.55 and 1.5);
\draw[curve](1.05,-.65)--(1.65,-.3)--(1.75,.65)--(1.15,.3)--cycle;
\draw[guide](.8,0)--(2.1,0);\draw[->](1.9,.3)arc(30:145:.45)node[above]{$\omega$};
\node[above]at(-1.3,1.6){$C_1,I_1$};\node[right]at(1.9,-.45){$R_2,L_2$};
\node[below]at(0,-1.9){$M=M_0\cos\omega t$};
\end{tikzpicture}
\caption{二次电流由感应决定，不再假定为外源维持的常量。}\label{fig:i2-h8-q}
\end{figure}
\begin{solution}
图~\ref{fig:i2-h8-q}中，一次电流固定，二次闭合回路的外磁链为$I_1M_0\cos\omega t$。感应电动势为$E_0\sin\omega t$，$E_0=I_1M_0\omega$。电路方程
\[L_2\dot I_2+R_2I_2=E_0\sin\omega t
\]
先取特解$I_{2p}=A\sin\omega t+B\cos\omega t$，其导数为$A\omega\cos\omega t-B\omega\sin\omega t$。比较正、余弦系数，
\[R_2A-\omega L_2B=E_0,\qquad \omega L_2A+R_2B=0,\]
第二式给$B=-\omega L_2A/R_2$，代入第一式，得$(R_2^2+\omega^2L_2^2)A=E_0R_2$。故$A=E_0R_2/(R_2^2+\omega^2L_2^2)$、$B=-E_0\omega L_2/(R_2^2+\omega^2L_2^2)$。齐次方程分离变量后有$\ln|I_{2h}|=-R_2t/L_2+C_0$，即$I_{2h}=Ce^{-R_2t/L_2}$。若初始电流为$I_{20}$，则$A\sin0+B\cos0+C=I_{20}$，所以$C=I_{20}-B$。题目求交流稳态，取衰减项消失后的特解，写成$I_2=(E_0/Z)\sin(\omega t-\delta)$，其中$Z=\sqrt{R_2^2+\omega^2L_2^2}$、$\tan\delta=\omega L_2/R_2$。图~\ref{fig:i2-h8-s}显示这一相位差。


\begin{figure}[H]\centering
\begin{tikzpicture}[>=Stealth,line cap=round,line join=round]
\draw[axis](0,0)--(4.4,0)node[right]{$\omega t$};\draw[axis](0,-1.2)--(0,1.5)node[above]{归一化量};
\draw[guide,domain=0:6.283,samples=101]plot({.6*\x},{sin(57.29578*\x)});
\draw[curve,domain=0:6.283,samples=101]plot({.6*\x},{sin(57.29578*\x-35)});
\node[above]at(2.4,1.3){电流滞后感应电动势};
\end{tikzpicture}
\caption{电感改变相位，从而改变平均机械阻力。}\label{fig:i2-h8-s}
\end{figure}

磁力矩按瞬时电流的共能导数为$\tau=I_1I_2\partial_\theta M=-I_1M_0I_2\sin\omega t$。周期平均时，
令$T=2\pi/\omega$，用积化和差公式积分，
\[\begin{aligned}
\langle\sin\omega t\sin(\omega t-\delta)\rangle
&=\frac1{2T}\int_0^T[\cos\delta-\cos(2\omega t-\delta)]\dd t\\
&=\frac1{2T}\left[t\cos\delta-\frac{\sin(2\omega t-\delta)}{2\omega}\right]_0^T
=\frac12\cos\delta.
\end{aligned}\]
因此
\[\boxed{\langle\tau\rangle
=-\frac{I_1^2M_0^2\omega R_2}{2(R_2^2+\omega^2L_2^2)}}.
\]
它反对给定转动，维持匀速的平均机械输入为$-\omega\langle\tau\rangle=E_0^2R_2/(2Z^2)=\langle I_2^2R_2\rangle$。

稳态二次磁能的平均变化率为零，固定$I_1$的源在一周期中也没有净交叉磁链输入；机械功最终成为焦耳热。若$R_2\to0$且仍有有限自感，平均阻力矩趋零，电流主要与驱动相差四分之一周期，能量在转动和磁场之间往返交换。不能把有自感的电流直接写成$\mathcal E/R_2$再取零电阻极限。
\end{solution}

\begin{exercise}[闭合无阻线圈的插入力]\label{ex:i2-h9}
两个无阻闭合线圈自感均为$L_0$，互感随插入量$x$变化为$M=cx$，$c>0$，且$M<L_0$。初始$x=0$，电流为$I_1=I_0,I_2=0$。求插入过程的电流与磁力，并用固定磁链能量核对。 见图~\ref{fig:i2-h9-q}。
\end{exercise}

\begin{figure}[H]\centering
\begin{tikzpicture}[>=Stealth,line cap=round,line join=round]
\draw[axis](-2.7,0)--(2.7,0)node[right]{$z$};
\draw[curve](-2.4,-.85)rectangle(.4,.85);
\foreach \t in{-2.25,-1.85,-1.45,-1.05,-.65,-.25,.15}{\draw[thin](\t,0)ellipse(.12 and .85);}
\draw[curve](0,-.43)rectangle(2.1,.43);
\foreach \t in{.15,.5,.85,1.2,1.55,1.9}{\draw[thin](\t,0)ellipse(.10 and .43);}
\node[above]at(-1.2,.95){$C_1,I_1$};\node[above]at(1.5,.55){$C_2,I_2$};
\draw[guide](0,-.5)--(0,-1.3);\draw[guide](.4,-.9)--(.4,-1.3);
\draw[<->](0,-1.18)--(.4,-1.18)node[midway,below]{$x$};
\draw[->](1.6,-.85)--(.8,-.85)node[midway,below]{插入};
\node[below]at(0,-1.8){两回路无阻，$L_1=L_2=L_0,\ M=cx$};
\end{tikzpicture}
\caption{忽略端场的线圈插入模型，$x$为重叠长度}\label{fig:i2-h9-q}
\end{figure}
\begin{solution}
图~\ref{fig:i2-h9-q}的两个回路都没有电源和电阻，所以保持的是各自磁链。初始$M=0,I_1=I_0,I_2=0$给出$\lambda_1=L_0I_0$、$\lambda_2=0$。插入后满足
\[L_0I_1+MI_2=L_0I_0,\qquad MI_1+L_0I_2=0.
\]
令$k=M/L_0<1$，第二式先给$I_2=-kI_1$，代回第一式得到$(1-k^2)I_1=I_0$。于是
\[\boxed{I_1=\frac{I_0}{1-k^2},\qquad
I_2=-\frac{kI_0}{1-k^2}}.
\]
副边产生反向感应电流。磁场力为$F_x=I_1I_2M'(x)$，故
\[\boxed{F_x=-\frac{ckI_0^2}{(1-k^2)^2}}.
\]
负号表示阻碍进一步插入。


\begin{figure}[H]\centering
\begin{tikzpicture}[>=Stealth,line cap=round,line join=round]
\draw[axis](-.1,0)--(4.5,0)node[right]{$k$};\draw[axis](0,-.1)--(0,2.5)node[above]{$W_B/(L_0I_0^2/2)$};
\draw[curve,domain=0:.83,samples=81]plot({4*\x},{1/(1-(\x)^2)});
\node[right]at(2.4,2){$1/(1-k^2)$};
\end{tikzpicture}
\caption{在固定磁链条件下，增大耦合需要增加磁能。}\label{fig:i2-h9-s}
\end{figure}

图~\ref{fig:i2-h9-s}的磁能为$W_B=(\lambda_1I_1+\lambda_2I_2)/2=L_0I_0^2/[2(1-k^2)]$。固定磁链求导时$k'=c/L_0$，故
\[-\frac{\dd W_B}{\dd x}=-\frac{L_0I_0^2}2
\frac{2kk'}{(1-k^2)^2}=-\frac{ckI_0^2}{(1-k^2)^2},\]
与共能法的瞬时力相同。准静态外力须输入正功以提高磁能。

若错误地固定$I_1=I_0,I_2=0$，便会得到零力，那描述的是另接电源强行取消感应电流的系统。$k\to1$时本题所需磁能发散，有限能量的自由运动无法达到该理想重合点；实际漏感和端场也会限制简单$M=cx$模型的适用范围。
\end{solution}

\begin{exercise}[交流恒电压下的气隙吸力]\label{ex:i2-h10}
线性铁芯绕$N$匝、截面$S$，有可变窄气隙$h$。绕组电阻忽略，接正弦端电压$\widehat V\cos\omega t$，磁通无直流偏置。求气隙瞬时与平均吸力，说明它为何不同于恒电流条件的气隙依赖。 见图~\ref{fig:i2-h10-q}。
\end{exercise}

\begin{figure}[H]\centering
\begin{tikzpicture}[>=Stealth,line cap=round,line join=round]
\fill[even odd rule,gray!18](-2,-1.4)rectangle(2,1.4)(-1.35,-.8)rectangle(1.35,.8);
\draw[curve](-2,-1.4)rectangle(2,1.4);\draw[curve](-1.35,-.8)rectangle(1.35,.8);
\foreach\y in{-.6,-.3,0,.3,.6}{\draw(-2.25,\y)--(-1.1,\y+.12);}
\node[left]at(-2.25,0){$N_1$};\draw[->](-.8,1.1)--(.8,1.1)node[midway,above]{$\Phi$};\fill[white](1.3,-.14)rectangle(2.1,.14);\node[right]at(2.2,0){$h$};\node[below]at(0,-1.9){$V_1=\widehat V\cos\omega t,\quad R_1=0$};
\end{tikzpicture}
\caption{电压积分规定磁通，磁阻规定励磁电流。}\label{fig:i2-h10-q}
\end{figure}
\begin{solution}
图~\ref{fig:i2-h10-q}给定绕组端电压，忽略电阻时$N\dot\Phi=\widehat V\cos\omega t$。取没有直流偏置的周期解，
\[\Phi(t)-\Phi(0)=\frac{\widehat V}{N}\int_0^t\cos\omega s\dd s
=\frac{\widehat V}{N}\left[\frac{\sin\omega s}{\omega}\right]_0^t.\]
无偏置解取$\Phi(0)=0$，所以
\[\Phi(t)=\frac{\widehat V}{N\omega}\sin\omega t,\qquad
B(t)=\frac{\widehat V}{NS\omega}\sin\omega t.
\]
磁通的幅值由电压积分决定，在理想条件下不随气隙改变；磁阻改变的是达到这条磁通所需的励磁电流。


\begin{figure}[H]\centering
\begin{tikzpicture}[>=Stealth,line cap=round,line join=round]
\draw[axis](0,0)--(4.4,0)node[right]{$\omega t$};\draw[axis](0,0)--(0,2.1)node[above]{$F$};
\draw[curve,domain=0:6.283,samples=101]plot({.6*\x},{1.7*(sin(57.29578*\x))^2});\draw[guide](0,.85)--(4,.85)node[right]{$\langle F\rangle$};
\end{tikzpicture}
\caption{吸力随磁通平方变化，具有非零平均值与两倍频分量。}\label{fig:i2-h10-s}
\end{figure}

由图~\ref{fig:i2-h10-s}所示的平方关系，气隙吸力为
\[F(t)=\frac{S B(t)^2}{2\mu_0}
=\frac{\widehat V^2}{2\mu_0N^2S\omega^2}\sin^2\omega t,
\]
所以
\[\boxed{\langle F\rangle=\frac{\widehat V^2}{4\mu_0N^2S\omega^2}}.
\]
这里用了$\sin^2\omega t=(1-\cos2\omega t)/2$，其中余弦项在整周期的积分为
\[\left[\frac{\sin2\omega t}{2\omega}\right]_0^{2\pi/\omega}=0.\]
平均力不含$h$，且有两倍频脉动。它并不与定电流时$F\propto(h+\ell_{\rm Fe}/\mu_r)^{-2}$矛盾，两题控制的是不同量。

当气隙增加，$I=\mathcal R_m\Phi/N$也增加；维持端电压的源负责相应的能量交换。若绕组电阻不可忽略，实际作用于磁链的电压为$V_1-RI$，气隙变化又通过电流改变这个电压，结论不再严格。题中还要求磁通幅值处于铁芯线性范围，且$h$小于截面尺度，使均匀气隙场近似有效。
\end{solution}

\begin{exercise}[变压器的励磁导纳]\label{ex:i2-h11}
共同磁通变压器匝数$N_1,N_2$，铁芯磁阻$\mathcal R_m$，二次接纯电阻$R_L$。采用$\ee^{\mathrm i\omega t}$和有效值相量，求一次输入导纳，画出电流相位关系，忽略绕组电阻和铁耗。 见图~\ref{fig:i2-h11-q}。
\end{exercise}

\begin{figure}[H]\centering
\begin{tikzpicture}[>=Stealth,line cap=round,line join=round]
\fill[even odd rule,gray!18](-2,-1.4)rectangle(2,1.4)(-1.35,-.8)rectangle(1.35,.8);
\draw[curve](-2,-1.4)rectangle(2,1.4);\draw[curve](-1.35,-.8)rectangle(1.35,.8);
\foreach\y in{-.6,-.3,0,.3,.6}{\draw(-2.25,\y)--(-1.1,\y+.12);}
\node[left]at(-2.25,0){$N_1$};\draw[->](-.8,1.1)--(.8,1.1)node[midway,above]{$\Phi$};\foreach\y in{-.5,-.25,0,.25,.5}{\draw(1.1,\y)--(2.25,\y+.1);}\node[right]at(2.25,0){$N_2,R_L$};\node[below]at(0,-1.9){有限磁阻$\mathcal R_m$，忽略铁耗};
\end{tikzpicture}
\caption{有限磁阻允许存在不供给负载的励磁电流。}\label{fig:i2-h11-q}
\end{figure}
\begin{solution}
图~\ref{fig:i2-h11-q}仍有共同磁通关系$\underline V_2=(N_2/N_1)\underline V_1$。二次电流为$\underline I_2=\underline V_2/R_L$，由安匝方程
\[\underline I_1=\frac{N_2}{N_1}\underline I_2+
\frac{\mathcal R_m}{N_1}\underline\Phi,\qquad
\underline V_1=\mathrm i\omega N_1\underline\Phi.
\]
因此
\[\boxed{Y_{\rm in}=\frac{\underline I_1}{\underline V_1}
=\frac{(N_2/N_1)^2}{R_L}+
\frac1{\mathrm i\omega L_m}},\qquad L_m=\frac{N_1^2}{\mathcal R_m}.
\]
第一项等于折算电阻$R_{\rm ref}=(N_1/N_2)^2R_L$的导纳，第二项为励磁电感导纳，两者并联。


\begin{figure}[H]\centering
\begin{tikzpicture}[>=Stealth,line cap=round,line join=round]
\draw[axis](0,0)--(3,0)node[right]{$\underline V_1$};\draw[->,thick](0,0)--(2,0)node[above]{$\underline I_{\rm load}$};\draw[->](2,0)--(2,-1.3)node[right]{$\underline I_m$};\draw[->,thick](0,0)--(2,-1.3)node[midway,below left]{$\underline I_1$};\node[below]at(0,-1.9){相量以有效值表示};
\end{tikzpicture}
\caption{负载电流与励磁电流分别位于相量的实部与负虚部方向。}\label{fig:i2-h11-s}
\end{figure}

图~\ref{fig:i2-h11-s}中，负载反射电流与电压同相，励磁电流滞后四分之一周期，总电流为两相量之和。若有效值电压为$V_1$，有功输入为$V_1^2/R_{\rm ref}$，无功为$V_1^2/(\omega L_m)$。这里采用有效值，不能再额外乘一个$1/2$。

若以一次电压相位为零，令$G=1/R_{\rm ref}$、$B_m=1/(\omega L_m)$，则
\[
\underline I_1=V_1(G-\mathrm iB_m),\qquad
|\underline I_1|=V_1\sqrt{G^2+B_m^2},\qquad
\arg\underline I_1=-\arctan\frac{B_m}{G}.
\]
图中水平向右的负载电流与竖直向下的励磁电流构成直角三角形，总电流滞后电压。复功率为
$\underline S=\underline V_1\underline I_1^{\,*}=V_1^2(G+\mathrm iB_m)$，实部、虚部分别就是前述有功和无功。

当$\mathcal R_m\to0$时，$L_m\to\infty$，励磁支路电流趋零，恢复纯负载折算。若再有一次串联电阻$R_1$，线圈端电压需由$\underline V_1=\underline U/(1+R_1Y_{\rm in})$求得，不能把电源电压直接当作绕组上的电压。
\end{solution}

\begin{exercise}[共同磁通三绕组的数值计算]\label{ex:i2-h12}
共同磁通理想变压器$N_1=600,N_2=200,N_3=100$，一次串联$R_1=2\Omega$，二次电阻$R_2=10\Omega,R_3=5\Omega$，电源$u=120\cos\omega t$伏。求各电流、电压及平均功率分配。 见图~\ref{fig:i2-h12-q}。
\end{exercise}

\begin{figure}[H]\centering
\begin{tikzpicture}[>=Stealth,line cap=round,line join=round]
\fill[even odd rule,gray!18](-2,-1.4)rectangle(2,1.4)(-1.35,-.8)rectangle(1.35,.8);
\draw[curve](-2,-1.4)rectangle(2,1.4);\draw[curve](-1.35,-.8)rectangle(1.35,.8);
\foreach\y in{-.6,-.3,0,.3,.6}{\draw(-2.25,\y)--(-1.1,\y+.12);}
\node[left]at(-2.25,0){$N_1$};\draw[->](-.8,1.1)--(.8,1.1)node[midway,above]{$\Phi$};\foreach\y in{.3,.5,.7}{\draw(1.1,\y)--(2.25,\y+.12);}\foreach\y in{-.7,-.5,-.3}{\draw(1.1,\y)--(2.25,\y+.12);}\node[right]at(2.5,.5){$N_2=200,\ R_2=10\Omega$};\node[right]at(2.5,-.4){$N_3=100,\ R_3=5\Omega$};\node[below]at(0,-1.9){$N_1=600,\ R_1=2\Omega,\ u=120\cos\omega t$};
\end{tikzpicture}
\caption{电源幅值为120伏，区别于有效值120伏。}\label{fig:i2-h12-q}
\end{figure}
\begin{solution}
图~\ref{fig:i2-h12-q}的两个负载都链接同一磁通，先计算$G_N=N_2^2/R_2+N_3^2/R_3=6000$。折算负载为$N_1^2/G_N=60\Omega$，与一次串联电阻相加为$62\Omega$，如图~\ref{fig:i2-h12-s}。


\begin{figure}[H]\centering
\begin{tikzpicture}[>=Stealth,line cap=round,line join=round]
\draw[curve](-2.4,0)--(2.4,0);\node[draw,fill=white]at(-1.2,0){$2\Omega$};\node[draw,fill=white]at(1,0){$60\Omega$};
\node[above]at(0,.9){$I_{1,\max}=60/31\ {\rm A}$};\node[below]at(0,-.9){线圈端电压低于电源电压};
\end{tikzpicture}
\caption{两个负载并联折算后，一次总阻值为62欧。}\label{fig:i2-h12-s}
\end{figure}

因此一次电流及一次线圈电压为
\[I_1=\frac{60}{31}\cos\omega t\ {\rm A},\qquad
V_1=\frac{3600}{31}\cos\omega t\ {\rm V}.
\]
按匝数比分配，$V_2=(1200/31)\cos\omega t$伏、$V_3=(600/31)\cos\omega t$伏；再各除以负载电阻，得到
\[\boxed{I_2=I_3=\frac{120}{31}\cos\omega t\ {\rm A}}.
\]
本题两个二次电流相等，是题给电阻与匝数共同作用的结果。改变任一负载后，共同磁通仍在，两电流一般不再相等。安匝校验为$600(60/31)=200(120/31)+100(120/31)$。

电源平均功率为$120(60/31)/2=3600/31$瓦。各电阻使用自己的电流幅值，分别得到
\[\begin{aligned}
\langle P_{R_1}\rangle&=\frac12\left(\frac{60}{31}\right)^2\!\times2=\frac{3600}{961}\ {\rm W},\\
\langle P_{R_2}\rangle&=\frac12\left(\frac{120}{31}\right)^2\!\times10=\frac{72000}{961}\ {\rm W},\\
\langle P_{R_3}\rangle&=\frac12\left(\frac{120}{31}\right)^2\!\times5=\frac{36000}{961}\ {\rm W}.
\end{aligned}\]
相加为$111600/961=3600/31$瓦，与电源输入相等。这里使用幅值，周期平均后有$1/2$；若混入上一题有效值相量的公式，功率会差一倍。
\end{solution}

\begin{exercise}[相同匝数但不同磁通结构]\label{ex:i2-h13}
双窗口分磁通理想铁芯中，一次和两个二次绕组的匝数均为$N$，两个二次分别接$R,3R$。一次链接两返回磁通之和。求输入电阻、电流和磁通幅值比，再与共同磁通三绕组结构比较。 见图~\ref{fig:i2-h13-q}。
\end{exercise}

\begin{figure}[H]\centering
\begin{tikzpicture}[>=Stealth,line cap=round,line join=round]
\draw[curve](-2,-1.2)rectangle(2,1.2);\draw[curve](-1.45,-.75)rectangle(-.3,.75);\draw[curve](.3,-.75)rectangle(1.45,.75);\node at(0,0){$N$};\node[left]at(-2.1,0){$N,R$};\node[right]at(2.1,0){$N,3R$};\draw[->](-.8,.95)--(-1.2,.95)node[above]{$\Phi_2$};\draw[->](.8,.95)--(1.2,.95)node[above]{$\Phi_3$};
\end{tikzpicture}
\caption{中柱绕组链接两窗口的磁通之和。}\label{fig:i2-h13-q}
\end{figure}
\begin{solution}
图~\ref{fig:i2-h13-q}有两个独立磁通窗口。两条理想磁路的安匝平衡分别给$NI_1=NI_2$、$NI_1=NI_3$，所以负载电流都等于一次电流。

负载电压满足$V_2=RI_1$、$V_3=3RI_1$，而一次链接两分支之和，故$V_1=N(\dot\Phi_2+\dot\Phi_3)=V_2+V_3$。由此
\[\boxed{R_{\rm in}=4R,\qquad I_2=I_3=I_1,\qquad V_3=3V_2}.
\]

\begin{figure}[H]\centering
\begin{tikzpicture}[>=Stealth,line cap=round,line join=round]
\draw[curve](-2.2,0)--(2.2,0);\node[draw,fill=white]at(-.9,0){$R$};\node[draw,fill=white]at(.9,0){$3R$};\node[above]at(0,.9){$I_2=I_3=I_1$};\node[below]at(0,-.9){$V_3=3V_2,\quad\dot\Phi_3=3\dot\Phi_2$};
\end{tikzpicture}
\caption{分磁通结构的各负载电流由各自安匝平衡确定。}\label{fig:i2-h13-s}
\end{figure}

图~\ref{fig:i2-h13-s}中折算电阻串联，与共同磁通并联结构不同。

在同频、无直流偏置的稳态，$\dot\Phi_3/\dot\Phi_2=V_3/V_2=3$，两个磁通幅值也为$3:1$。负载功率比为$P_3:P_2=3:1$，总功率$4RI_1^2$等于一次输入。

若把相同的三个匝数放到一个共同磁通窗口上，反而会得到$V_2=V_3$、$I_2=3I_3$及并联负载$R\parallel3R=3R/4$。电气负载与匝数可以完全相同，铁芯所链接的磁通不同，端口等效关系就不同。计算前必须从图上识别每个绕组围住的磁通，不能只凭匝数套用同一个变压器公式。
\end{solution}

\begin{exercise}[多分支磁路的方程数]\label{ex:i2-h14}
一条带$N$匝励磁线圈的公共磁路，磁阻$\mathcal R_c$，与三条磁阻分别为$\mathcal R,2\mathcal R,3\mathcal R$的返回支路组成连通网络。求磁通分配和线圈电感，并说明应列多少条独立节点与回路方程。 见图~\ref{fig:i2-h14-q}。
\end{exercise}

\begin{figure}[H]\centering
\begin{tikzpicture}[>=Stealth,line cap=round,line join=round]
\coordinate(A)at(0,1.6);\coordinate(B)at(0,-1.6);\foreach\x in{-2,-.7,.7,2}{\draw[curve](A)--(\x,1)--(\x,-1)--(B);}
\node[left]at(-2,0){$NI,\mathcal R_c$};\node[draw,fill=white]at(-.7,0){$\mathcal R$};\node[draw,fill=white]at(.7,0){$2\mathcal R$};\node[draw,fill=white]at(2,0){$3\mathcal R$};
\fill(A)circle(1.4pt)node[above]{$A$};\fill(B)circle(1.4pt)node[below]{$B$};
\end{tikzpicture}
\caption{四支路连接同一对磁节点，公共支路含励磁安匝。}\label{fig:i2-h14-q}
\end{figure}
\begin{solution}
图~\ref{fig:i2-h14-q}的两节点之间有公共励磁支路和三条返回支路。若无漏磁，主磁通$\Phi$等于三个返回磁通之和。返回支路承受同一磁动势降，故磁通比为倒磁阻比$1:1/2:1/3=6:3:2$。

图~\ref{fig:i2-h14-s}中三者分别为$6\Phi/11,3\Phi/11,2\Phi/11$。并联磁阻满足$\mathcal R_p^{-1}=1/\mathcal R+1/(2\mathcal R)+1/(3\mathcal R)=11/(6\mathcal R)$，于是
\[\boxed{\mathcal R_{\rm eq}=\mathcal R_c+\frac{6\mathcal R}{11}},\qquad
\boxed{L=\frac{N^2}{\mathcal R_c+6\mathcal R/11}}.
\]

\begin{figure}[H]\centering
\begin{tikzpicture}[>=Stealth,line cap=round,line join=round]
\draw[curve](0,0)rectangle(4.4,1);\draw[curve](2.4,0)--(2.4,1);\draw[curve](3.6,0)--(3.6,1);\fill[gray!18](0,0)rectangle(2.4,1);\fill[gray!10](2.4,0)rectangle(3.6,1);
\node at(1.2,.5){$6/11$};\node at(3,.5){$3/11$};\node at(4,.5){$2/11$};\node[below]at(2.2,-.5){三分支磁通占主磁通的比例};
\end{tikzpicture}
\caption{磁通按各分支磁导分配，总和为主磁通。}\label{fig:i2-h14-s}
\end{figure}

该图有四条支路、两个节点，节点磁通连续方程只有一条独立；还需$4-2+1=3$条独立Ampere回路方程。三个返回支路分别与励磁支路组成回路，就给出一组可用的方程。

不能把三个返回支路各自当作一个独立线圈，再把其电感普通并联，因为它们共享励磁安匝和公共支路磁阻。若公共磁阻增大，全部主磁通减少，但在返回磁阻比不变的近似下，图中的三种分配比例仍不变。
\end{solution}

